\section{Derivaciones y diferenciales}
\label{section:0.20-es}
\label{0.20-es}

Las nociones introducidas en este parágrafo se retomarán en forma
geométrica en el capítulo~IV, §~16, y desempeñarán un papel importante en el
estudio de los preesquemas. Su importancia en este capítulo se debe, en primer
lugar, a sus relaciones con la noción de álgebra formalmente lisa y, en
particular, a los teoremas (20.4.9), (20.5.7) y (20.5.12), que se traducirán al
lenguaje geométrico en el parágrafo del capítulo~IV dedicado a los morfismos
lisos y son de uso constante en las aplicaciones. Por otra parte, las nociones
diferenciales servirán para demostrar en el §~22 importantes criterios de
regularidad que desempeñarán un papel esencial en el estudio profundo de los
anillos locales noetherianos realizado en el §~7 del capítulo~IV.

\subsection{Derivaciones y extensiones de álgebras}
\label{subsection:0.20.1-es}

\oldpage[0\textsubscript{IV-1}]{117}
\begin{proposition}[20.1.1]
\label{0.20.1.1-es}
Sean $A$ un anillo (no necesariamente conmutativo), $B$, $E$, $C$
$A$-anillos, $p:E\to C$ un $A$-homomorfismo cuyo núcleo $\mathfrak{Q}$
es de cuadrado nulo, $u:B\to C$ un $A$-homomorfismo. Supongamos que existe
un $A$-homomorfismo $v_0:B\to E$ tal que $u$ se factoriza como
$B\xrightarrow{v_0}E\xrightarrow{p}C$. Entonces el conjunto de los
$A$-homomorfismos $v:B\to E$ tales que $u=p\circ v$ es idéntico al
conjunto de las aplicaciones $v_0+D$, donde $D:B\to\mathfrak{Q}$ es una
aplicación que verifica las dos condiciones siguientes:
\begin{enumerate}
  \item[\rm(i)] $D$ es un homomorfismo de $A$-bimódulos.
  \item[\rm(ii)] Cualesquiera que sean $f$, $g$ en $B$, se tiene
\end{enumerate}
\begin{equation}
\tag{20.1.1.1}
\label{0.20.1.1.1-es}
  D(fg)=f.D(g)+D(f).g.
\end{equation}
\end{proposition}

\begin{proof}
Decir que $p(v(f))=p(v_0(f))$ para $f\in B$ significa que
$D(f)=v(f)-v_0(f)$ pertenece a $\mathfrak{Q}$; al escribir que
$v(fg)=v(f)v(g)$, se obtiene la relación \sref{0.20.1.1.1-es},
puesto que $\mathfrak{Q}$ es de cuadrado nulo, y la condición (i) resulta de que
$v$ y $v_0$ son $A$-homomorfismos.

Si $\rho:A\to B$ es el homomorfismo estructural, se deduce de
\sref{0.20.1.1.1-es} que se tiene
$D(\rho(a)f)=\rho(a)D(f)+D(\rho(a))f$ para todo $a\in A$; pero se debe
tener $D(\rho(a)f)=\rho(a)D(f)$ por (i), luego, tomando $f=1$, se
obtiene
\[
  D(\rho(a))=0 \qquad\text{para}\quad a\in A\,;
\]
recíprocamente, si $D$ es nula en $\rho(A)$ y verifica
\sref{0.20.1.1.1-es}, también verifica (i).
\end{proof}

\begin{definition}[20.1.2]
\label{0.20.1.2-es}
Dados un $A$-anillo $B$ y un $B$-bimódulo $L$, se llama
\emph{$A$-derivación de $B$ en $L$} a una aplicación $D:B\to L$ que verifica
las condiciones (i) y (ii) de \sref{0.20.1.1-es}.
\end{definition}

Se deduce de \sref{0.20.1.1-es} que el \emph{núcleo} de una $A$-derivación
$D:B\to L$ es un sub-$A$-anillo de $B$.

A veces se llama \emph{derivación} de $B$ en $L$ a una
$\mathbf{Z}$-derivación, es decir, una aplicación \emph{aditiva} de $B$
en $L$ que verifica \sref{0.20.1.1.1-es}; se tiene, por tanto, $D(1)=0$ para tal
aplicación. Si $B$ es un álgebra sobre un \emph{cuerpo primo} $P$, toda
$\mathbf{Z}$-derivación de $B$ es una $P$-derivación: esto es evidente por lo
anterior si $P$ es de característica $>0$ y, en caso contrario,
la relación $n.n^{-1}=1$ para $n\in\mathbf{Z}$ da
\[
  nD(n^{-1})=0
\]
luego $D(n^{-1})=0$, puesto que $P$ es de característica $0$.

Se deduce inmediatamente de la definición \sref{0.20.1.2-es} que, si $D$, $D'$
son dos $A$-derivaciones de $B$ en $L$, lo mismo ocurre con $D-D'$. En otras
palabras, el conjunto de las $A$-derivaciones de $B$ en $L$ está provisto de una
estructura de \emph{grupo aditivo}; este grupo se denota por
$\operatorname{D\acute{e}r}_A(B,L)$. Si $A$ es \emph{conmutativo}, $B$ un
$A$-\emph{álgebra conmutativa}, y $L$ un $B$-módulo, entonces, para toda
$A$-derivación $D$ de $B$ en $L$ y todo $a\in A$, $aD$ sigue siendo una
$A$-derivación de $B$ en $L$; en otras palabras, $\operatorname{D\acute{e}r}_A(B,L)$
está entonces provisto de una estructura de $A$-\emph{módulo}.

\oldpage[0\textsubscript{IV-1}]{118}
La prop. \sref{0.20.1.1-es} se interpreta de la manera siguiente:

\begin{corollary}[20.1.3]
\label{0.20.1.3-es}
Dados dos $A$-homomorfismos de $A$-anillos $p:E\to C$,
$u:B\to C$, el primero de los cuales tiene un núcleo $\mathfrak{Q}$ de cuadrado
nulo, el conjunto de los $A$-homomorfismos $v:B\to E$ tales que $u=p\circ v$
es vacío o es un espacio homogéneo principal bajo el grupo
$\operatorname{D\acute{e}r}_A(B,\mathfrak{Q})$.
\end{corollary}

En particular:

\begin{corollary}[20.1.4]
\label{0.20.1.4-es}
Sean $A$ un anillo, $C$ un $A$-anillo, $L$ un $C$-bimódulo, $E$ una
$A$-extensión de $C$ por $L$, $p:E\to C$ la aumentación. La aplicación que
asocia a toda derivación $D\in\operatorname{D\acute{e}r}_A(C,L)$ la aplicación
$x\mapsto x+D(p(x))$ es un isomorfismo del grupo
$\operatorname{D\acute{e}r}_A(C,L)$ sobre el grupo de las $A$-equivalencias de $E$
sobre sí misma.
\end{corollary}

\begin{proof}
Apliquemos \sref{0.20.1.1-es} para $B=E$, $v_0=\mathrm{I}_E$: el conjunto
de los $A$-homomorfismos $v:E\to E$ tales que $p\circ v=p$ es idéntico al
conjunto de las aplicaciones $v=\mathrm{I}_E+D'$, donde
$D'\in\operatorname{D\acute{e}r}_A(E,L)$. Decir que tal $A$-homomorfismo $v$ es
una $A$-equivalencia equivale a decir que la restricción de $v$ a $L$ es la
inyección canónica, o también que $D'(x)=0$ en $L$; pero esto significa asimismo
que $D'$ se factoriza como $E\xrightarrow{p}C\xrightarrow{D}L$, donde $D$ es una
$A$-derivación. De donde resulta el corolario.
\end{proof}

Para las extensiones triviales, \sref{0.20.1.3-es} da:

\begin{corollary}[20.1.5]
\label{0.20.1.5-es}
Sean $A$ un anillo, $B$, $C$ dos $A$-anillos, $L$ un $C$-bimódulo,
$u:B\to C$ un $A$-homomorfismo; la aplicación que asocia a toda derivación
$D\in\operatorname{D\acute{e}r}_A(B,L)$ la aplicación
$\alpha:x\mapsto(u(x),D(x))$ es una biyección sobre el conjunto de los
$A$-homomorfismos $v:B\to D_C(L)$ (cf. 18.2.3) tales que $u$ se factoriza como
$B\xrightarrow{v}D_C(L)\to C$.
\end{corollary}

Más particularmente:

\begin{corollary}[20.1.6]
\label{0.20.1.6-es}
Sean $A$ un anillo, $B$ un $A$-anillo, $L$ un $B$-bimódulo. Si a toda
derivación $D\in\operatorname{D\acute{e}r}_A(B,L)$ se asocian: $1^\circ$ la
$A$-equivalencia $(x,y)\mapsto(x,y+D(x))$ de la extensión $D_B(L)$ sobre
sí misma; $2^\circ$ el $A$-homomorfismo $x\mapsto(x,D(x))$ de $B$ en
$D_B(L)$, inverso por la derecha del homomorfismo de aumentación $D_B(L)\to B$,
se definen correspondencias biunívocas canónicas entre:
\begin{enumerate}
  \item[\rm(i)] el conjunto $\operatorname{D\acute{e}r}_A(B,L)$;
  \item[\rm(ii)] el conjunto de las $A$-equivalencias de $D_B(L)$ sobre sí misma;
  \item[\rm(iii)] el conjunto de los $A$-homomorfismos inversos por la derecha del
  homomorfismo de aumentación $D_B(L)\to B$.
\end{enumerate}
\end{corollary}

Se observará que la correspondencia biunívoca así establecida entre (i) y (ii)
respeta las estructuras de grupo, y que la que se deduce de ella entre (ii) y
(iii) no es sino la correspondencia biunívoca ya definida en (18.3.8).

\begin{corollary}[20.1.7]
\label{0.20.1.7-es}
Sean $A$ un anillo, $C$ un $A$-anillo, $L$ un $A$-bimódulo, $E$ una
$A$-extensión trivial de $C$ por $L$, $p:E\to C$ la aumentación. El conjunto
de las $A$-derivaciones $D$ de $E$ en $L$ tales que $D(x)=x$ en $L$ es
idéntico al conjunto de las aplicaciones $x\mapsto x-w(p(x))$, donde $w$
recorre el conjunto de los $A$-homomorfismos inversos por la derecha de $p$.
\end{corollary}

\begin{proof}
Basta aplicar de nuevo \sref{0.20.1.1-es} para $B=E$,
$v_0=\mathrm{I}_E$; si $v=\mathrm{I}_E-D$, la condición $D(x)=x$ para
$x\in L$ equivale a $v(x)=0$ para $x\in L$, es decir, a $v=w\circ p$,
donde $w:C\to E$ es un $A$-homomorfismo; además, la condición
$p\circ v=p$ equivale a $p\circ w=\mathrm{I}_C$, en otras palabras, al hecho de
que $w$ sea inverso por la derecha de $p$.
\end{proof}

\subsection{Propiedades funtoriales de las derivaciones}
\label{subsection:0.20.2-es}

\oldpage[0\textsubscript{IV-1}]{119}
\begin{env}[20.2.1]
\label{0.20.2.1-es}
Sean $A$ un anillo, $B$ un $A$-anillo, $L$ un $B$-bimódulo; si $L'$ es
un segundo $B$-bimódulo y $w:L\to L'$ un homomorfismo de $B$-bimódulos, es
claro que la aplicación $D\mapsto w\circ D$ es un homomorfismo de
grupos aditivos
\begin{equation}
\tag{20.2.1.1}
\label{0.20.2.1.1-es}
  w_0:\operatorname{D\acute{e}r}_A(B,L)\to\operatorname{D\acute{e}r}_A(B,L')
\end{equation}
y que, si $w':L'\to L''$ es un segundo homomorfismo de $B$-bimódulos, se tiene
$(w'\circ w)_0=w'_0\circ w_0$. Cuando $A$ es conmutativo, $B$ un
$A$-álgebra conmutativa y $L$ un $B$-módulo, \sref{0.20.2.1.1-es} es un
homomorfismo de $A$-módulos.

En segundo lugar, sean $B'$ un $A$-anillo, $v:B'\to B$ un
$A$-homomorfismo que hace de $L$ un $B'$-bimódulo; entonces la aplicación
$D\mapsto D\circ v$ es un homomorfismo de grupos aditivos,
\begin{equation}
\tag{20.2.1.2}
\label{0.20.2.1.2-es}
  v^0:\operatorname{D\acute{e}r}_A(B,L)\to\operatorname{D\acute{e}r}_A(B',L)
\end{equation}
como se deduce de \sref{0.20.1.1.1-es}; si $v':B''\to B'$ es un segundo
$A$-homomorfismo, se tiene $(v\circ v')^0=v'^0\circ v^0$. Cuando $A$, $B$ y
$B'$ son conmutativos y $L$ un $B$-módulo, \sref{0.20.2.1.2-es} es un
homomorfismo de $A$-módulos.

Finalmente, sea $u:A'\to A$ un homomorfismo de anillos que hace de $B$ un
$A'$-anillo; toda $A$-derivación $D\in\operatorname{D\acute{e}r}_A(B,L)$ es también
una $A'$-derivación, de donde resulta una inyección canónica de grupos conmutativos
\begin{equation}
\tag{20.2.1.3}
\label{0.20.2.1.3-es}
  u^0:\operatorname{D\acute{e}r}_A(B,L)\to\operatorname{D\acute{e}r}_{A'}(B,L)
\end{equation}
y, si $u':A''\to A'$ es un segundo homomorfismo de anillos, se tiene
$(u\circ u')^0=u'^0\circ u^0$; cuando $A$, $A'$ y $B$ son conmutativos y
$L$ un $B$-módulo, \sref{0.20.2.1.3-es} es un dihomomorfismo de módulos
(relativo a $u$).

También se puede decir que
\[
  (A,B,L)\leadsto\operatorname{D\acute{e}r}_A(B,L)
\]
es un \emph{funtor covariante} de la categoría $\boldsymbol{K}$ definida en
(18.3.5) en la categoría $\boldsymbol{Ab}$ de los grupos conmutativos, haciendo
corresponder a toda terna $(u,v,w)$ que constituye un morfismo de
$\boldsymbol{K}$ el homomorfismo $w_0\circ v^0\circ u^0$; la verificación de
la funtorialidad resulta de la conmutatividad de los diagramas
\[
\xymatrix{
  \operatorname{D\acute{e}r}_A(B,L) \ar[r]^{v^0} \ar[d]_{w_0} &
    \operatorname{D\acute{e}r}_A(B',L) \ar[d]^{w_0} \\
  \operatorname{D\acute{e}r}_A(B,L') \ar[r]^{v^0} &
    \operatorname{D\acute{e}r}_A(B',L')
}
\qquad
\xymatrix{
  \operatorname{D\acute{e}r}_A(B,L) \ar[r]^{u^0} \ar[d]_{w_0} &
    \operatorname{D\acute{e}r}_{A'}(B,L) \ar[d]^{w_0} \\
  \operatorname{D\acute{e}r}_A(B,L') \ar[r]^{u^0} &
    \operatorname{D\acute{e}r}_{A'}(B,L')
}
\]
para todo homomorfismo $w:L\to L'$ de $B$-bimódulos.
\end{env}

\oldpage[0\textsubscript{IV-1}]{120}
\begin{theorem}[20.2.2]
\label{0.20.2.2-es}
Sean $u:A\to B$, $v:B\to C$ dos homomorfismos de anillos, $L$ un
$C$-bimódulo. Se tiene una sucesión exacta canónica de grupos conmutativos
\begin{align}
\tag{20.2.2.1}
\label{0.20.2.2.1-es}
  0\to\operatorname{D\acute{e}r}_B(C,L)\xrightarrow{u^0}
  \operatorname{D\acute{e}r}_A(C,L)&\xrightarrow{v^0}
  \operatorname{D\acute{e}r}_A(B,L)\xrightarrow{\partial} \\
  &\to\operatorname{Exan}_B(C,L)\xrightarrow{u^1}
  \operatorname{Exan}_A(C,L)\xrightarrow{v^1}
  \operatorname{Exan}_A(B,L) \notag
\end{align}
donde $u^0$, $v^0$ son los homomorfismos (20.2.1.3) y (20.2.1.2)
respectivamente, $u^1$, $v^1$ los homomorfismos definidos en (18.3.4.1) y
(18.3.3.1) respectivamente, y donde $\partial$ se define como sigue: para
toda $A$-derivación $D$ de $B$ en $L$, $\partial(D)$ es la clase de la
$B$-extensión de $C$ por $L$ definida sobre el anillo $D_C(L)$ por el
$A$-homomorfismo $\alpha:x\mapsto(v(x),D(x))$ (cf.
\sref{0.20.1.5-es}). Además, la sucesión exacta \sref{0.20.2.2.1-es} es
funtorial en $L$ (para los homomorfismos definidos en (20.2.1.1) y
(18.3.1.1) respectivamente).
\end{theorem}

\begin{proof}
Como $D$ es una $A$-derivación (y \emph{a fortiori} una
$\mathbf{Z}$-derivación) de $B$ en $L$, el $A$-homomorfismo
$\alpha:x\mapsto(v(x),D(x))$ define efectivamente sobre $D_C(L)$ una estructura de
$B$-extensión, por tanto $\partial(D)$ está bien definida \sref{0.20.1.5-es}. Hay
que verificar la exactitud en $5$ lugares:

\medskip
\noindent 1)\enspace
La exactitud en $\operatorname{D\acute{e}r}_B(C,L)$ es trivial (cf.
\sref{0.20.2.1-es}).

\medskip
\noindent 2)\enspace
Por definición \sref{0.20.2.1-es} el núcleo de $v^0$ es el conjunto de las
$A$-derivaciones de $C$ en $L$ que se anulan en $v(B)$, es decir,
las $A$-derivaciones que son también $B$-derivaciones
\sref{0.20.1.1-es}; de donde resulta la exactitud en $\operatorname{D\acute{e}r}_A(C,L)$.

\medskip
\noindent 3)\enspace
El núcleo de $\partial$ está formado por las derivaciones
$D\in\operatorname{D\acute{e}r}_A(B,L)$ para las cuales la $B$-extensión definida por
$\alpha:x\mapsto(v(x),D(x))$ es $B$-trivial; esto significa (18.2.3) que
existe un $B$-homomorfismo $z\mapsto(z,w(z))$ de $C$ en $D_C(L)$ (estando la
estructura de $B$-anillo de $D_C(L)$ definida por $\alpha$); pero un tal
homomorfismo, al ser \emph{a fortiori} un $A$-homomorfismo, es de la forma
$z\mapsto(z,D'(z))$ donde $D'\in\operatorname{D\acute{e}r}_A(C,L)$
\sref{0.20.1.6-es}; y al escribir que es un $B$-homomorfismo, se obtiene
\[
  D(x)z+v(x)D'(z)=D'(v(x)z)=v(x)D'(z)+D'(v(x))z
\]
para $x\in B$, $z\in C$, lo que da $D'(v(x))=D(x)$; el núcleo de
$\partial$ es por tanto la imagen de $v^0$.

\medskip
\noindent 4)\enspace
El núcleo de $u^1$ es el conjunto de las clases de $B$-extensiones de $C$ por
$L$ que son $A$-triviales (18.3.7), por tanto (salvo equivalencia) de la forma
$D_C(L)$, donde la estructura de $A$-anillo está definida por el homomorfismo
$t\mapsto(u(t),0)$. Ahora bien, toda estructura de $B$-extensión sobre $D_C(L)$ está
definida por un homomorfismo $\alpha:x\mapsto(v(x),D(x))$, donde $D$ es una
$\mathbf{Z}$-derivación de $B$ en $L$ \sref{0.20.1.5-es}; decir que la
estructura de $A$-anillo de esta $B$-extensión se deduce de su estructura de
$B$-anillo por medio de $u:A\to B$ significa que $D(u(t))=0$ para $t\in A$,
por tanto que $D$ es una $A$-derivación o también que la clase de la $B$-extensión
considerada es de la forma $\partial(D)$; de donde resulta la exactitud en
$\operatorname{Exan}_B(C,L)$.

\medskip
\noindent 5)\enspace
El núcleo de $v^1$ es el conjunto de las clases de las $A$-extensiones $E$ de $C$
por $L$ que se trivializan sobre $v(B)$, es decir, para las cuales existe un
$A$-homomorfismo $w:B\to E$ tal que $v$ se factoriza como
$B\xrightarrow{w}E\to C$; ahora bien, un tal $A$-homomorfismo define sobre $E$ una
estructura de $B$-extensión cuya clase tiene por imagen mediante $u^1$ la clase de
la $A$-extensión dada; siendo trivial el recíproco, se ha probado
la exactitud en $\operatorname{Exan}_A(C,L)$.

Finalmente, la funtorialidad en $L$ se deduce trivialmente de las definiciones.
\end{proof}

\oldpage[0\textsubscript{IV-1}]{121}
\begin{corollary}[20.2.3]
\label{0.20.2.3-es}
Sean $A$, $B$, $C$ tres anillos conmutativos, $u:A\to B$, $v:B\to C$
dos homomorfismos de anillos, $L$ un $C$-módulo. Se tiene una sucesión exacta
canónica de $A$-módulos
\begin{align}
\tag{20.2.3.1}
\label{0.20.2.3.1-es}
  0\to\operatorname{D\acute{e}r}_B(C,L)\xrightarrow{u^0}
  \operatorname{D\acute{e}r}_A(C,L)&\xrightarrow{v^0}
  \operatorname{D\acute{e}r}_A(B,L)\xrightarrow{\partial} \\
  &\to\operatorname{Exalcom}_B(C,L)\xrightarrow{u^1}
  \operatorname{Exalcom}_A(C,L)\xrightarrow{v^1}
  \operatorname{Exalcom}_A(B,L) \notag
\end{align}
funtorial en $L$.
\end{corollary}

\begin{proof}
El razonamiento es el mismo que en \sref{0.20.2.2-es} una vez que se ha
verificado que para toda derivación $D\in\operatorname{D\acute{e}r}_A(B,L)$,
$\partial(D)$ es efectivamente la clase de una $B$-extensión de $C$ por $L$ que es
una $B$-álgebra conmutativa; pero esto resulta inmediatamente de la conmutatividad
de $C$ y del hecho de que $L$ es un $C$-módulo.
\end{proof}

\begin{corollary}[20.2.4]
\label{0.20.2.4-es}
Bajo las hipótesis de \sref{0.20.2.2-es} (resp.
\sref{0.20.2.3-es}), se tiene una sucesión exacta canónica, funtorial en $L$,
\begin{equation}
\tag{20.2.4.1}
\label{0.20.2.4.1-es}
  0\to\operatorname{D\acute{e}r}_B(C,L)\xrightarrow{u^0}
  \operatorname{D\acute{e}r}_A(C,L)\xrightarrow{v^0}
  \operatorname{D\acute{e}r}_A(B,L)\xrightarrow{\partial}
  \operatorname{Exan}_{B/A}(C,L)\to0
\end{equation}
(resp.
\begin{equation}
\tag{20.2.4.2}
\label{0.20.2.4.2-es}
  0\to\operatorname{D\acute{e}r}_B(C,L)\xrightarrow{u^0}
  \operatorname{D\acute{e}r}_A(C,L)\xrightarrow{v^0}
  \operatorname{D\acute{e}r}_A(B,L)\xrightarrow{\partial}
  \operatorname{Exalcom}_{B/A}(C,L)\to0\text{).}
\end{equation}
Esto resulta de la definición de $\operatorname{Exan}_{B/A}(C,L)$ (resp.
$\operatorname{Exalcom}_{B/A}(C,L)$) ((18.3.7) y (18.4.2)).
\end{corollary}

\begin{remark}[20.2.5]
\label{0.20.2.5-es}
Supongamos que se tiene un diagrama conmutativo de homomorfismos de anillos
\[
\xymatrix{
  A \ar[r] & B \ar[r] & C \\
  A' \ar[u] \ar[r] & B' \ar[u] \ar[r] & C' \ar[u]
}
\]
Entonces se tiene un diagrama conmutativo
\[
\xymatrix@C=0.5em{
  0 \ar[r] & \operatorname{D\acute{e}r}_B(C,L) \ar[r] \ar[d] &
  \operatorname{D\acute{e}r}_A(C,L) \ar[r] \ar[d] &
  \operatorname{D\acute{e}r}_A(B,L) \ar[r]^-\partial \ar[d] &
  \operatorname{Exan}_B(C,L) \ar[r] \ar[d] &
  \operatorname{Exan}_A(C,L) \ar[r] \ar[d] &
  \operatorname{Exan}_A(B,L) \ar[d] \\
  0 \ar[r] & \operatorname{D\acute{e}r}_{B'}(C',L) \ar[r] &
  \operatorname{D\acute{e}r}_{A'}(C',L) \ar[r] &
  \operatorname{D\acute{e}r}_{A'}(B',L) \ar[r]^-\partial &
  \operatorname{Exan}_{B'}(C',L) \ar[r] &
  \operatorname{Exan}_{A'}(C',L) \ar[r] &
  \operatorname{Exan}_{A'}(B',L)
}
\]
y del mismo modo para las sucesiones exactas \sref{0.20.2.3.1-es},
\sref{0.20.2.4.1-es} y \sref{0.20.2.4.2-es}.
\end{remark}

\subsection{Derivaciones continuas en los anillos topológicos}
\label{subsection:0.20.3-es}

\begin{env}[20.3.1]
\label{0.20.3.1-es}
Dados dos anillos topológicos $A$, $B$ (linealmente topologizados
como siempre), denotaremos por $\operatorname{Hom.\ cont}(A,B)$
el conjunto de los homomorfismos continuos de $A$ en $B$. Dado un anillo
topológico $A$, la categoría de los $A$-anillos topológicos se define como
la de los $A$-anillos (18.1.4), sustituyendo en todas partes «anillo» por «anillo
topológico» y «homomorfismo» por «homomorfismo continuo»; si $B$

\oldpage[0\textsubscript{IV-1}]{122}
y $C$ son dos $A$-anillos topológicos, se denotará por
$\operatorname{Hom.\ cont}_A(B,C)$ el conjunto de los $A$-homomorfismos continuos
de $B$ en $C$.

Sean $A$ un anillo topológico, $B$ un $A$-anillo topológico, $L$ un
$B$-bimódulo topológico; se denota por $\operatorname{D\acute{e}r.\ cont}_A(B,L)$
el conjunto de las $A$-derivaciones \emph{continuas} de $B$ en $L$; está claro
que es un \emph{subgrupo} de $\operatorname{D\acute{e}r}_A(B,L)$ (y un
sub-$A$-módulo cuando $A$ y $B$ son conmutativos y $L$ un $B$-módulo).
\end{env}

\begin{env}[20.3.2]
\label{0.20.3.2-es}
Es inmediato que la prop. \sref{0.20.1.1-es} subsiste cuando se sustituye en ella
«anillo» por «anillo topológico» y «homomorfismo» por
«homomorfismo continuo». Del mismo modo, \sref{0.20.1.3-es} y
\sref{0.20.1.4-es} siguen siendo válidos sustituyendo $\operatorname{D\acute{e}r}$ por
$\operatorname{D\acute{e}r.\ cont}$, «anillo» (resp. «bimódulo») por «anillo
topológico» (resp. «bimódulo topológico»), «homomorfismo» por
«homomorfismo continuo»; naturalmente hay que suponer en
\sref{0.20.1.4-es} que $p$ es continuo, y sustituir «$A$-equivalencias» por
«$A$-equivalencias continuas». Se tienen los mismos resultados para
\sref{0.20.1.5-es} y \sref{0.20.1.6-es}, a condición de tomar como
topología sobre $D_C(L)$ (resp. $D_B(L)$) la topología \emph{producto} (sobre
$C\times L$, resp. $B\times L$).
\end{env}

\begin{proposition}[20.3.3]
\label{0.20.3.3-es}
Sean $A$ un anillo topológico, $B$ un $A$-anillo topológico, $L$ un
$B$-bimódulo topológico discreto anulado por un ideal bilateral abierto de $B$.
Si en $B$ el cuadrado de todo ideal bilateral abierto es abierto, se tiene
$\operatorname{D\acute{e}r.\ cont}_A(B,L)=\operatorname{D\acute{e}r}_A(B,L)$.
\end{proposition}

\begin{proof}
En efecto, si $\mathfrak{R}$ es un ideal bilateral abierto de $B$ que anula a $L$,
y $D$ una $A$-derivación de $B$ en $L$, se tiene
$D(\mathfrak{R}^2)\subset\mathfrak{R}.D(\mathfrak{R})+
D(\mathfrak{R}).\mathfrak{R}=0$ \sref{0.20.1.1.1-es}, por tanto $D$ es continua.
\end{proof}

\begin{env}[20.3.4]
\label{0.20.3.4-es}
Todos los resultados de \sref{0.20.2.1-es} siguen siendo válidos cuando se sustituye
«anillo» por «anillo topológico», «bimódulo» por «bimódulo
topológico», «homomorfismo» por «homomorfismo continuo» y
$\operatorname{D\acute{e}r}$ por $\operatorname{D\acute{e}r.\ cont}$.
\end{env}

\begin{proposition}[20.3.5]
\label{0.20.3.5-es}
Sean $A$ un anillo topológico, $B$ un $A$-anillo topológico, $L$ un
$B$-bimódulo topológico discreto anulado por un ideal bilateral abierto de $B$.
Se tiene entonces un isomorfismo canónico
\begin{equation}
\tag{20.3.5.1}
\label{0.20.3.5.1-es}
  \varinjlim\operatorname{D\acute{e}r}_{A/\mathfrak{J}}
  (B/\mathfrak{R},L)\xrightarrow{\sim}
  \operatorname{D\acute{e}r.\ cont}_A(B,L)
\end{equation}
donde en el primer miembro el límite inductivo se toma según el conjunto
ordenado filtrante de los pares $(\mathfrak{J},\mathfrak{R})$ de ideales bilaterales
tales que $\mathfrak{R}.L=L.\mathfrak{R}=0$,
$\mathfrak{J}.B\subset\mathfrak{R}$,
$B.\mathfrak{J}\subset\mathfrak{R}$.
\end{proposition}

\begin{proof}
Como $A/\mathfrak{J}$ y $B/\mathfrak{R}$ son discretos, se tienen
homomorfismos canónicos
$w_{\mathfrak{J},\mathfrak{R}}:
\operatorname{D\acute{e}r}_{A/\mathfrak{J}}(B/\mathfrak{R},L)\to
\operatorname{D\acute{e}r.\ cont}_A(B,L)$ que forman un sistema inductivo
\sref{0.20.3.4-es}, de donde se obtiene el homomorfismo \sref{0.20.3.5.1-es} por paso al
límite inductivo. Como el homomorfismo
$B/\mathfrak{R}'\to B/\mathfrak{R}$ es sobreyectivo para
$\mathfrak{R}\supset\mathfrak{R}'$, resulta inmediatamente de la definición que
el homomorfismo
$\operatorname{D\acute{e}r}_{A/\mathfrak{J}}(B/\mathfrak{R},L)\to
\operatorname{D\acute{e}r}_{A/\mathfrak{J}}(B/\mathfrak{R}',L)$ (con
$\mathfrak{R}.L=L.\mathfrak{R}=0$,
$\mathfrak{J}.B\subset\mathfrak{R}'$,
$B.\mathfrak{J}\subset\mathfrak{R}'$) es inyectivo, y evidentemente ocurre
lo mismo con el homomorfismo
$\operatorname{D\acute{e}r}_{A/\mathfrak{J}}(B/\mathfrak{R},L)\to
\operatorname{D\acute{e}r}_{A/\mathfrak{J}'}(B/\mathfrak{R},L)$ para
$\mathfrak{J}'\subset\mathfrak{J}$ \sref{0.20.2.1-es}; se concluye que
el homomorfismo \sref{0.20.3.5.1-es} es inyectivo. Por otra parte, si $D$ es
una $A$-derivación continua de $B$ en $L$, su núcleo contiene un ideal
bilateral abierto $\mathfrak{R}_0$ de $B$, y si $\mathfrak{J}_0$ es un ideal
bilateral abierto de $A$ tal que $\mathfrak{J}_0B\subset\mathfrak{R}_0$ y
$B\mathfrak{J}_0\subset\mathfrak{R}_0$, está claro que $D$ es la imagen
canónica de una $(A/\mathfrak{J}_0)$-derivación de
$B/\mathfrak{R}_0$ en $L$, luego \sref{0.20.3.5.1-es} es sobreyectivo.
\end{proof}

\oldpage[0\textsubscript{IV-1}]{123}
\begin{proposition}[20.3.6]
\label{0.20.3.6-es}
Sean $u:A\to B$, $v:B\to C$ dos homomorfismos continuos de anillos
topológicos, $L$ un $C$-bimódulo discreto anulado por un ideal bilateral abierto
de $C$. Se tiene una sucesión exacta canónica
\begin{align}
\tag{20.3.6.1}
\label{0.20.3.6.1-es}
  0\to\operatorname{D\acute{e}r.\ cont}_B(C,L)\xrightarrow{u^0}
  \operatorname{D\acute{e}r.\ cont}_A(C,L)&\xrightarrow{v^0}
  \operatorname{D\acute{e}r.\ cont}_A(B,L)\xrightarrow{\partial} \\
  &\to\operatorname{Exantop}_B(C,L)\xrightarrow{u^1}
  \operatorname{Exantop}_A(C,L)\xrightarrow{v^1}
  \operatorname{Exantop}_A(B,L) \notag
\end{align}
donde $\partial$ se define por paso al límite inductivo a partir del
homomorfismo $\partial$ de \sref{0.20.2.2.1-es}; esta sucesión exacta es
funtorial en $L$ (en la categoría de los $C$-bimódulos discretos anulados por
ideales bilaterales abiertos).
\end{proposition}

\begin{proof}
Esto resulta de la exactitud del funtor $\varinjlim$, a partir de
\sref{0.20.2.2-es}.
\end{proof}

\begin{corollary}[20.3.7]
\label{0.20.3.7-es}
Sean $A$, $B$, $C$ tres anillos topológicos conmutativos, $u:A\to B$,
$v:B\to C$ dos homomorfismos continuos, $L$ un $C$-módulo discreto anulado
por un ideal abierto de $C$. Se tiene una sucesión exacta canónica de $A$-módulos,
funtorial en $L$,
\begin{align}
\tag{20.3.7.1}
\label{0.20.3.7.1-es}
  &0\to\operatorname{D\acute{e}r.\ cont}_B(C,L)\xrightarrow{u^0}
  \operatorname{D\acute{e}r.\ cont}_A(C,L)\xrightarrow{v^0}
  \operatorname{D\acute{e}r.\ cont}_A(B,L)\xrightarrow{\partial} \\
  &\to\operatorname{Exalcotop}_B(C,L)\xrightarrow{u^1}
  \operatorname{Exalcotop}_A(C,L)\xrightarrow{v^1}
  \operatorname{Exalcotop}_A(B,L). \notag
\end{align}
\end{corollary}

\begin{corollary}[20.3.8]
\label{0.20.3.8-es}
Bajo las hipótesis de \sref{0.20.3.5-es} (resp.
\sref{0.20.3.6-es}) se tiene una sucesión exacta canónica, funtorial en $L$
\begin{equation}
\tag{20.3.8.1}
\label{0.20.3.8.1-es}
  0\to\operatorname{D\acute{e}r.\ cont}_B(C,L)\xrightarrow{u^0}
  \operatorname{D\acute{e}r.\ cont}_A(C,L)\xrightarrow{v^0}
  \operatorname{D\acute{e}r.\ cont}_A(B,L)\to
  \operatorname{Exantop}_{B/A}(C,L)\to0
\end{equation}
(resp.
\begin{equation}
\tag{20.3.8.2}
\label{0.20.3.8.2-es}
  0\to\operatorname{D\acute{e}r.\ cont}_B(C,L)\xrightarrow{u^0}
  \operatorname{D\acute{e}r.\ cont}_A(C,L)\xrightarrow{v^0}
  \operatorname{D\acute{e}r.\ cont}_A(B,L)\to
  \operatorname{Exalcotop}_{B/A}(C,L)\to0\text{).}
\end{equation}
\end{corollary}

Dejamos al lector la tarea de escribir los diagramas análogos a los de
\sref{0.20.2.5-es}.

\subsection{Partes principales y diferenciales}
\label{subsection:0.20.4-es}

En todo lo que sigue de este parágrafo y en los tres siguientes, todos los
anillos se suponen \emph{conmutativos}.

\begin{env}[20.4.1]
\label{0.20.4.1-es}
Sean $A$ un anillo topológico, $B$ una $A$-álgebra topológica; la
$A$-álgebra $B\otimes_A B$ estará provista de la topología \emph{producto
tensorial}, que la convierte en una $A$-álgebra topológica; denotaremos por $p$
(o $p_{B/A}$) el $A$-homomorfismo canónico sobreyectivo
\begin{equation}
\tag{20.4.1.1}
\label{0.20.4.1.1-es}
  p:B\otimes_A B\to B
\end{equation}
tal que $p(b\otimes b')=bb'$; es inmediato que $p$ es continuo. El
\emph{núcleo} de $p$ se denotará por $\mathfrak{J}_{B/A}$ (o solamente
$\mathfrak{J}$ si no hay confusión). Denotaremos por
$j_1:B\to B\otimes_A B$ y $j_2:B\to B\otimes_A B$ los dos
$A$-homomorfismos canónicos, tales que
\[
  j_1(b)=b\otimes1, \qquad j_2(b)=1\otimes b
\]
que son continuos.
\end{env}

\oldpage[0\textsubscript{IV-1}]{124}
\begin{definition}[20.4.2]
\label{0.20.4.2-es}
Se llama \emph{$B$-álgebra aumentada de partes principales de orden $1$ de
$B$ respecto de $A$} y se denota por $\mathrm{P}^1_{B/A}$ la $A$-álgebra
topológica cociente
\begin{equation}
\tag{20.4.2.1}
\label{0.20.4.2.1-es}
  \mathrm{P}^1_{B/A}=(B\otimes_A B)/\mathfrak{J}^2
\end{equation}
provista de la estructura de $B$-álgebra topológica definida por el homomorfismo
$\bar{j}_1:B\to\mathrm{P}^1_{B/A}$ (deducido de $j_1$ por composición con
el homomorfismo canónico $B\otimes_A B\to\mathrm{P}^1_{B/A}$), y de
la aumentación de $B$-álgebra $\varepsilon:\mathrm{P}^1_{B/A}\to B$ (también
denotada por $\varepsilon_{B/A}$) deducida de $p$ por paso al cociente.

Como $p(b\otimes1)=b$ por definición, está claro que $\varepsilon$ es
efectivamente una aumentación de $B$-álgebra.
\end{definition}

\begin{definition}[20.4.3]
\label{0.20.4.3-es}
El núcleo de la aumentación $\varepsilon:\mathrm{P}^1_{B/A}\to B$,
\begin{equation}
\tag{20.4.3.1}
\label{0.20.4.3.1-es}
  \Omega^1_{B/A}=\mathfrak{J}_{B/A}/(\mathfrak{J}_{B/A})^2
\end{equation}
provisto de la topología inducida por la de $\mathrm{P}^1_{B/A}$, que lo convierte
en un $B$-módulo topológico, se llama el $B$-módulo de las
\emph{$1$-diferenciales} (o simplemente de las \emph{diferenciales}) de $B$
respecto de $A$.
\end{definition}

Obsérvese que la topología de $\Omega^1_{B/A}$ es también la topología cociente
de la topología inducida sobre $\mathfrak{J}_{B/A}$ por la de
$B\otimes_A B$ (Bourbaki, \emph{Top. gén.}, cap.~III, 3\textsuperscript{a}
ed., §~2, n\textsuperscript{o}~7, prop.~20). Si $B$ es discreto, lo mismo ocurre
con $\Omega^1_{B/A}$. Se denota por $\widehat{\Omega}^1_{B/A}$ el
\emph{completado separado} del $B$-módulo topológico $\Omega^1_{B/A}$.

Todo anillo topológico $B$ puede considerarse como una
$\mathbf{Z}$-álgebra topológica (donde $\mathbf{Z}$ está provisto de la topología
discreta), por lo que se puede definir el $B$-módulo topológico $\Omega^1_{B/\mathbf{Z}}$,
que también se llama a veces el $B$-módulo de las \emph{diferenciales absolutas}
sobre $B$ y se denota por $\Omega^1_B$. Si $B$ es un álgebra topológica sobre
un cuerpo primo $P$ (discreto), se tiene
$B\otimes_P B=B\otimes_{\mathbf{Z}}B$, luego
$\Omega^1_{B/P}=\Omega^1_B$.

\begin{lemma}[20.4.4]
\label{0.20.4.4-es}
Sean $A$ un anillo, $B$ una $A$-álgebra. El ideal $\mathfrak{J}_{B/A}$ de
$B\otimes_A B$ está generado por los elementos $1\otimes s-s\otimes1$, donde $s$
recorre un conjunto de generadores de la $A$-álgebra $B$.
\end{lemma}

\begin{proof}
Está claro que, para todo $x\in B$, se tiene
$x\otimes1-1\otimes x\in\mathfrak{J}$; por otra parte, cualesquiera que sean $x$,
$y$ en $B$, se tiene
$x\otimes y=xy\otimes1+(x\otimes1)(1\otimes y-y\otimes1)$. Si
$\sum_i(x_i\otimes y_i)\in\mathfrak{J}$, se tiene por definición
$\sum_i x_i y_i=0$, luego
\begin{equation}
\tag{20.4.4.1}
\label{0.20.4.4.1-es}
  \sum_i(x_i\otimes y_i)=
  \sum_i(x_i\otimes1)(1\otimes y_i-y_i\otimes1)
\end{equation}
lo que demuestra que $\mathfrak{J}$ es el ideal generado por los elementos
$1\otimes x-x\otimes1$. Además, si $x=st$, se tiene
\begin{equation}
\tag{20.4.4.2}
\label{0.20.4.4.2-es}
  x\otimes1-1\otimes x=(s\otimes1)(t\otimes1-1\otimes t)
  +(s\otimes1-1\otimes s)(1\otimes t)
\end{equation}
lo que concluye inmediatamente la demostración por recurrencia.
\end{proof}

\begin{proposition}[20.4.5]
\label{0.20.4.5-es}
Sean $A$ un anillo topológico, $B$ una $A$-álgebra topológica. La
topología de $\Omega^1_{B/A}$ es menos fina que la topología deducida de la
de $B$ (19.0.2); si en $B$ el cuadrado de todo ideal abierto es abierto, estas
dos topologías son idénticas.
\end{proposition}

\begin{proof}
La primera aserción es trivial, pues la topología producto tensorial sobre
$B\otimes_A B$ es menos fina que la topología deducida de la de $B$;
\emph{a fortiori}, la topología inducida sobre
$\mathfrak{J}=\mathfrak{J}_{B/A}$

\oldpage[0\textsubscript{IV-1}]{125}
por la de $B\otimes_A B$ es menos fina que la topología sobre $\mathfrak{J}$
deducida de la de $B$. Para demostrar la segunda aserción, escribamos
Denotemos por $M\otimes N$, por abuso de notación, el submódulo
$\operatorname{Im}(M\otimes_A N)$, para dos sub-$A$-módulos $M$, $N$ de $B$.
Utilizando la relación
\[
\begin{aligned}
  (xy)\otimes z-x\otimes(yz)
  &=(x\otimes1)(1\otimes z)(y\otimes1-1\otimes y)\\
  &=xz.(y\otimes1-1\otimes y)
    -x.(z\otimes1-1\otimes z)(y\otimes1-1\otimes y)
\end{aligned}
\]
en el $B$-módulo $B\otimes_A B$ (definido por $j_1$), se ve inmediatamente,
teniendo en cuenta \sref{0.20.4.4-es}, que, si $\mathfrak{R}$ es un ideal de
$B$, se tiene
\[
  ((\mathfrak{R}^2\otimes B)+(B\otimes\mathfrak{R}^2))\cap\mathfrak{J}
  \subset(\mathfrak{R}\otimes\mathfrak{R})\cap\mathfrak{J}
  +\mathfrak{R}\mathfrak{J}+\mathfrak{J}^2
\]
y por otra parte, si $x_i$, $y_i$ son elementos de $\mathfrak{R}$ tales
que $\sum_i(x_i\otimes y_i)\in\mathfrak{J}$, se deduce de
\sref{0.20.4.4.1-es} que se tiene
$\sum_i(x_i\otimes y_i)\in\mathfrak{R}\mathfrak{J}$, de modo que finalmente
\begin{equation}
\tag{20.4.5.1}
\label{0.20.4.5.1-es}
  (\mathfrak{R}^2\otimes B+B\otimes\mathfrak{R}^2)\cap\mathfrak{J}
  \subset\mathfrak{R}\mathfrak{J}+\mathfrak{J}^2.
\end{equation}
Ahora bien, se obtiene un sistema fundamental de entornos de $0$ en $\mathfrak{J}$
(para la topología inducida por la de $B\otimes_A B$) tomando como
entornos de $0$ los conjuntos
$(\mathfrak{R}\otimes B+B\otimes\mathfrak{R})\cap\mathfrak{J}$, donde
$\mathfrak{R}$ recorre el conjunto de los ideales abiertos. Como la topología de
$\Omega^1_{B/A}$ deducida de la de $B$ es también el cociente por
$\mathfrak{J}^2$ de la topología de $\mathfrak{J}$ deducida de la de $B$,
la hipótesis sobre los ideales abiertos de $B$ y la relación
\sref{0.20.4.5.1-es} terminan de demostrar la segunda aserción.
\end{proof}

\begin{definition}[20.4.6]
\label{0.20.4.6-es}
Sean $p_1$ y $p_2$ las aplicaciones compuestas
$B\xrightarrow{j_1}B\otimes_A B\to\mathrm{P}^1_{B/A}$,
$B\xrightarrow{j_2}B\otimes_A B\to\mathrm{P}^1_{B/A}$, que son
$A$-homomorfismos continuos tales que
$\varepsilon\circ p_1=\varepsilon\circ p_2=\mathrm{I}_B$. El
$A$-homomorfismo continuo de $A$-módulos
\begin{equation}
\tag{20.4.6.1}
\label{0.20.4.6.1-es}
  d_{B/A}=p_1-p_2:B\to\Omega^1_{B/A}
\end{equation}
se denomina la \emph{diferencial exterior de $B$ relativa a $A$}; para
todo $x\in B$, $d_{B/A}(x)$ (denotado también $d(x)$ o $dx$) se denomina la
\emph{diferencial de $x$} (relativa a $A$).
\end{definition}

Cuando $A=\mathbf{Z}$, se escribe $d_B$ en lugar de $d_{B/\mathbf{Z}}$; si $B$
es un álgebra sobre un cuerpo primo $P$, se tiene $d_B=d_{B/P}$.

\begin{proposition}[20.4.7]
\label{0.20.4.7-es}
El $B$-módulo $\Omega^1_{B/A}$ está generado por los elementos $d_{B/A}(x)$,
donde $x$ recorre un sistema de generadores del $A$-álgebra $B$.
\end{proposition}

\begin{proof}
Como $d_{B/A}(x)$ es la imagen canónica de $x\otimes1-1\otimes x$ en
$\Omega^1_{B/A}$, la proposición es una consecuencia inmediata de
\sref{0.20.4.4-es}.
\end{proof}

\begin{theorem}[20.4.8]
\label{0.20.4.8-es}
Sean $A$ un anillo topológico y $B$ un $A$-álgebra topológica.
\begin{enumerate}
  \item[\rm(i)] Existe un único isomorfismo de $B$-álgebras topológicas
  aumentadas
\begin{equation}
\tag{20.4.8.1}
\label{0.20.4.8.1-es}
  \varphi:\mathrm{P}^1_{B/A}\xrightarrow{\sim}D_B(\Omega^1_{B/A})
\end{equation}
cuya restricción a $\Omega^1_{B/A}$ es la identidad.
  \item[\rm(ii)] El homomorfismo $d_{B/A}$ es una $A$-derivación de $B$
  en $\Omega^1_{B/A}$, con la siguiente propiedad universal: para todo
  $B$-módulo topológico $L$, la aplicación $u\mapsto u\circ d_{B/A}$ es
  un isomorfismo de $A$-módulos
\begin{equation}
\tag{20.4.8.2}
\label{0.20.4.8.2-es}
  \operatorname{Hom.\ cont}_B(\Omega^1_{B/A},L)
  \xrightarrow{\sim}\operatorname{D\acute{e}r.\ cont}_A(B,L).
\end{equation}
\end{enumerate}
\end{theorem}

\begin{proof}
\oldpage[0\textsubscript{IV-1}]{126}
(i) Es inmediato que $\varphi$ (con las notaciones de
\sref{0.20.4.6-es}) es necesariamente la aplicación
$z\mapsto(\varepsilon(z),z-p_1(\varepsilon(z)))$, y el isomorfismo
inverso la aplicación $(b,x)\mapsto p_1(b)+x$; como estas dos aplicaciones
son continuas, esto demuestra la primera aserción. Obsérvese que esto
implica que la topología de $B$ se identifica (mediante $p_1$) con el \emph{cociente}
de la topología de $B\otimes_A B$ por $\mathfrak{J}_{B/A}$.

(ii) El hecho de que $d_{B/A}$ sea una $A$-derivación de $B$ se deduce de la
definición \sref{0.20.4.6-es} y de \sref{0.20.1.1-es}.

Para demostrar la propiedad universal, recordemos que
$\operatorname{D\acute{e}r.\ cont}_A(B,L)$ se identifica canónicamente con
el conjunto $G$ de los homomorfismos continuos de $A$-álgebras
$u:B\to D_B(L)$ tales que la composición
$B\xrightarrow{u}D_B(L)\xrightarrow{q}B$ sea la identidad (donde $q(b,0)=b$ es
la proyección; cf. \sref{0.20.1.6-es} y \sref{0.20.3.2-es}). Por otra parte,
gracias al isomorfismo $\varphi$,
$\operatorname{Hom.\ cont}_B(\Omega^1_{B/A},L)$ se identifica canónicamente con
el conjunto de los homomorfismos continuos de $B$-álgebras
$v:\mathrm{P}^1_{B/A}\to D_B(L)$ tales que la composición
$\mathrm{P}^1_{B/A}\xrightarrow{v}D_B(L)\xrightarrow{q}B$ sea
el aumento $\varepsilon$. Como por definición se tiene $p_2=p_1-d_{B/A}$,
basta demostrar que todo $u\in G$ se factoriza como
\[
\xymatrix{
  B \ar[r]^u \ar[d]_{p_2} & D_B(L) \\
  \mathrm{P}^1_{B/A} \ar[ur]_v &
}
\]
donde $v$ es un $B$-homomorfismo continuo. Ahora bien, ya se tiene un homomorfismo
continuo de $A$-álgebras $j:b\mapsto(b,0)$ de $B$ en $D_B(L)$, que
pertenece a $G$; por la definición del producto tensorial topológico
de álgebras topológicas ($\mathbf{0_I}$, 7.7.6), existe por tanto un
$A$-homomorfismo continuo de álgebras
$w:B\otimes_A B\to D_B(L)$ que hace conmutativo el diagrama
\[
\xymatrix{
  B\otimes_A B \ar[dr]^w & B \ar[l]_{j_1} \ar[d]^j \\
  B \ar[u]^{j_2} \ar[r]_u & D_B(L)
}
\]
Se tiene, pues, por definición
$w(b\otimes1-1\otimes b)=j(b)-u(b)\in L$, y en virtud de
\sref{0.20.4.4-es}, esto implica $w(\mathfrak{J})\subset L$ y por tanto
$w(\mathfrak{J}^2)=0$; en consecuencia, $w$ se factoriza como
\[
  B\otimes_A B\to\mathrm{P}^1_{B/A}\xrightarrow{v}D_B(L)
\]
donde $v$ es un homomorfismo continuo de $A$-álgebras; además, como
$v\circ p_1=j$ es un homomorfismo de $B$-álgebras, lo mismo ocurre con
$v$ por definición de la estructura de $B$-álgebra de
$\mathrm{P}^1_{B/A}$; como por definición se tiene $u=v\circ p_2$, esto
termina la demostración.
\end{proof}

\begin{theorem}[20.4.9]
\label{0.20.4.9-es}
Supongamos que $B$ es un $A$-álgebra topológica formalmente lisa. Entonces
el $B$-módulo topológico $\Omega^1_{B/A}$ es \emph{formalmente proyectivo}.
\end{theorem}

\begin{proof}
En efecto, la hipótesis implica que $B\otimes_A B$ (dotado de la estructura de
$B$-álgebra topológica definida por $j_1$) es un $B$-álgebra topológica
formalmente lisa \sref{0.19.3.5-es}, (iii), y por consiguiente también un
$A$-álgebra topológica formalmente lisa \sref{0.19.3.5-es}, (ii); como
$B$ es topológicamente isomorfa al $A$-álgebra cociente
$(B\otimes_A B)/\mathfrak{J}$ y es un $A$-álgebra formalmente lisa, la
conclusión se deduce de \sref{0.19.5.3-es}.
\end{proof}

\oldpage[0\textsubscript{IV-1}]{127}
\begin{corollary}[20.4.10]
\label{0.20.4.10-es}
Supongamos además que en $B$ el cuadrado de todo ideal abierto es abierto;
entonces, para todo ideal abierto $\mathfrak{R}$ de $B$,
$\Omega^1_{B/A}\otimes_B(B/\mathfrak{R})
=\Omega^1_{B/A}/\mathfrak{R}.\Omega^1_{B/A}$ es un
$(B/\mathfrak{R})$-módulo proyectivo.
\end{corollary}

\begin{proof}
En efecto, la topología de $\Omega^1_{B/A}$ es entonces la deducida de la de
$B$ \sref{0.20.4.5-es}, y basta aplicar \sref{0.19.2.4-es}.
\end{proof}

\begin{corollary}[20.4.11]
\label{0.20.4.11-es}
Sean $A$, $B$ dos anillos locales noetherianos y $\rho:A\to B$ un
homomorfismo local que hace de $B$ un $A$-álgebra topológica formalmente
lisa (para las topologías preádicas). Entonces, para todo ideal de definición
$\mathfrak{b}$ de $B$,
$\Omega^1_{B/A}\otimes_B(B/\mathfrak{b})$ es un
$(B/\mathfrak{b})$-módulo libre.
\end{corollary}

\begin{proof}
En efecto, de \sref{0.20.4.10-es} se deduce que este módulo es proyectivo, y
como $B/\mathfrak{b}$ es un anillo artiniano, todo
$(B/\mathfrak{b})$-módulo proyectivo es libre
($\mathbf{0_{III}}$, 10.1.3).
\end{proof}

\begin{proposition}[20.4.12]
\label{0.20.4.12-es}
Sean $A$ un anillo topológico y $B$ un $A$-álgebra topológica. Si
el homomorfismo estructural $A\to B$ es sobreyectivo, se tiene
$\Omega^1_{B/A}=0$.
\end{proposition}

\begin{proof}
En efecto, se tiene $\operatorname{D\acute{e}r}_A(B,L)=0$ para todo $B$-módulo
$L$ en virtud de \sref{0.20.1.1-es}, y la proposición se deduce inmediatamente
de \sref{0.20.4.8-es}.
\end{proof}

\begin{examples}[20.4.13]
\label{0.20.4.13-es}
\begin{enumerate}
  \item[\rm(i)] Sean $A$ un anillo y
  $B=A[X_\alpha]_{\alpha\in I}$ un álgebra de polinomios sobre $A$. Entonces
  $\Omega^1_{B/A}$ es un $B$-módulo libre, y los $dX_\alpha$ forman una
  base. En efecto, los $dX_\alpha$ generan este $B$-módulo
  \sref{0.20.4.7-es}. Por otra parte, si $L$ es un $B$-módulo libre, con
  una base $(z_\alpha)_{\alpha\in I}$ cuyo conjunto de índices es $I$, existe
  un $A$-homomorfismo $u$ de $B$ en $D_B(L)$ tal que
  $u(X_\alpha)=(X_\alpha,z_\alpha)$, para todo $\alpha\in I$, y por tanto, por
  \sref{0.20.1.5-es}, una $A$-derivación $D$ de $B$ en $L$ tal que
  $D(X_\alpha)=z_\alpha$ para todo $\alpha$; en virtud de
  \sref{0.20.4.8.1-es}, esto demuestra que los $dX_\alpha$ son linealmente
  independientes.
  \item[\rm(ii)] Sean $A$ un anillo, $L$ un $A$-módulo y $B$ la
  $A$-álgebra $D_A(L)$; entonces el homomorfismo canónico
  $x\mapsto d_{B/A}x$ de $L$ en $\Omega^1_{B/A}$ es \emph{biyectivo}, pues
  es inmediato que las $B$-derivaciones de $B=D_A(L)$ en un $B$-módulo
  $M$ son las aplicaciones de la forma $(b,x)\mapsto u(x)$, donde
  $u\in\operatorname{Hom}_A(L,M)$; se concluye, pues, por
  \sref{0.20.4.8-es}, (ii).
  \item[\rm(iii)] Supongamos que $B$ es el producto de dos $A$-álgebras
  topológicas $B_1$, $B_2$ (identificadas con ideales de $B$). Entonces las
  imágenes de $B_1\otimes_A B_2$ y de $B_2\otimes_A B_1$ por el homomorfismo
  $p$ \sref{0.20.4.1.1-es} son nulas, de donde se deduce inmediatamente que
  $\mathrm{P}^1_{B/A}$ se identifica con el \emph{producto}
  $\mathrm{P}^1_{B_1/A}\times\mathrm{P}^1_{B_2/A}$, y que el $B$-módulo
  $\Omega^1_{B/A}$ es \emph{suma directa} (topológica) de los $B$-módulos
  $\Omega^1_{B_1/A}$ y $\Omega^1_{B_2/A}$ (anulados respectivamente por
  $B_2$ y $B_1$).
  \item[\rm(iv)] Sean $A$, $B$ dos anillos íntegros tales que $A\subset
  B$, que $A$ sea integralmente cerrado, $B$ entero sobre $A$, y que el cuerpo
  de fracciones de $B$ sea una extensión separable del de $A$. Entonces
  el $B$-módulo $\Omega^1_{B/A}$ es un \emph{módulo de torsión}. En efecto,
  para todo $x\in B$, el polinomio minimal de $x$ respecto del cuerpo de
  fracciones $K$ de $A$ es un polinomio $f(T)$ perteneciente a $A[T]$; como
  $x$ es separable sobre $K$, se tiene $f'(x)\neq0$ y, por otra parte, de
  la relación $f(x)=0$ se deduce que $f'(x)d_{B/A}x=0$, de donde se sigue
  la aserción en virtud de \sref{0.20.4.7-es}.
\end{enumerate}
\end{examples}

\begin{remarks}[20.4.14]
\label{0.20.4.14-es}
\begin{enumerate}
  \item[\rm(i)] Obsérvese que el $B$-módulo $\Omega^1_{B/A}$, privado de su
  topología, es independiente de las topologías de $A$ y de $B$.
  \item[\rm(ii)] Introduciremos más adelante el «álgebra de partes
  principales de orden $n$» de $B$ respecto de $A$,
  $\mathrm{P}^n_{B/A}=(B\otimes_A B)/(\mathfrak{J}_{B/A})^{n+1}$, que constituye
  la base del «cálculo diferencial de orden $n$».
\end{enumerate}
\end{remarks}

\subsection{Propiedades funtoriales fundamentales de
\texorpdfstring{$\Omega^1_{B/A}$}{Omega1(B/A)}}
\label{subsection:0.20.5-es}

\oldpage[0\textsubscript{IV-1}]{128}
\begin{env}[20.5.1]
\label{0.20.5.1-es}
En todo este número y el siguiente, salvo mención expresa en contrario, se
supone que los anillos y módulos considerados están provistos de la topología
\emph{discreta}.
\end{env}

\begin{env}[20.5.2]
\label{0.20.5.2-es}
Sean $A$ un anillo, $B$, $C$ dos $A$-álgebras, $u:B\to C$ un
$A$-homomorfismo; se tiene un diagrama conmutativo
\[
\xymatrix{
  B\otimes_A B \ar[r]^{u\otimes u} \ar[d]_{p_{B/A}} &
  C\otimes_A C \ar[d]^{p_{C/A}} \\
  B \ar[r]^u & C
}
\tag{20.5.2.1}
\label{0.20.5.2.1-es}
\]
de donde se obtiene, pasando a los cocientes, un $A$-homomorfismo de álgebras
\begin{equation}
\tag{20.5.2.2}
\label{0.20.5.2.2-es}
  u':\mathrm{P}^1_{B/A}\to\mathrm{P}^1_{C/A}
\end{equation}
tal que el diagrama
\[
\xymatrix{
  \mathrm{P}^1_{B/A} \ar[r]^{u'} & \mathrm{P}^1_{C/A} \\
  B \ar[u]^{p_1} \ar[r]^u & C \ar[u]_{p_1}
}
\]
sea conmutativo; como $u\otimes u$ aplica $\mathfrak{J}_{B/A}$ en
$\mathfrak{J}_{C/A}$, se obtiene, restringiendo $u'$ a
$\Omega^1_{B/A}$, una aplicación
\begin{equation}
\tag{20.5.2.3}
\label{0.20.5.2.3-es}
  u'':\Omega^1_{B/A}\to\Omega^1_{C/A}
\end{equation}
tal que el par $(u'',u)$ sea un \emph{dihomomorfismo} para la
estructura de $B$-módulo de $\Omega^1_{B/A}$ y la estructura de $C$-módulo de
$\Omega^1_{C/A}$; este último hecho permite deducir de él canónicamente un
\emph{homomorfismo de $C$-módulos}
\begin{equation}
\tag{20.5.2.4}
\label{0.20.5.2.4-es}
  u_{C/B/A}:\Omega^1_{B/A}\otimes_B C\to\Omega^1_{C/A}.
\end{equation}

\oldpage[0\textsubscript{IV-1}]{129}
Además, como el diagrama
\[
\xymatrix{
  \mathrm{P}^1_{B/A} \ar[r]^{u'} & \mathrm{P}^1_{C/A} \\
  B \ar[u]^{p_2} \ar[r]^u & C \ar[u]_{p_2}
}
\tag{20.5.2.5}
\label{0.20.5.2.5-es}
\]
es también conmutativo, se deduce que el diagrama
\[
\xymatrix{
  \Omega^1_{B/A} \ar[r]^{u''} & \Omega^1_{C/A} \\
  B \ar[u]^{d_{B/A}} \ar[r]^u & C \ar[u]_{d_{C/A}}
}
\tag{20.5.2.6}
\label{0.20.5.2.6-es}
\]
es conmutativo.

Por último, si $w:C\to D$ es un segundo homomorfismo de $A$-álgebras, se tiene la
propiedad de transitividad
\begin{equation}
\tag{20.5.2.7}
\label{0.20.5.2.7-es}
  (w\circ u)_{D/B/A}
  =w_{D/C/A}\circ(u_{C/B/A}\otimes\mathrm{I}_D)
\end{equation}
como resulta de la definición.
\end{env}

\begin{env}[20.5.3]
\label{0.20.5.3-es}
Sean ahora $A$, $B$ dos anillos, $v:A\to B$ un homomorfismo
de anillos, $C$ una $B$-álgebra que se convierte en una $A$-álgebra mediante $v$;
entonces la aplicación canónica $v_0:C\otimes_A C\to C\otimes_B C$ es un
\emph{dihomomorfismo} sobreyectivo de álgebras (relativo a $v:A\to B$) tal que
el diagrama
\[
\xymatrix{
  C\otimes_A C \ar[r]^{v_0} \ar[d]_{p_{C/A}} &
  C\otimes_B C \ar[d]^{p_{C/B}} \\
  C \ar[r]^{\mathrm{I}_C} & C
}
\tag{20.5.3.1}
\label{0.20.5.3.1-es}
\]
sea conmutativo; pasando a los cocientes, se deduce un
dihomomorfismo de álgebras
\begin{equation}
\tag{20.5.3.2}
\label{0.20.5.3.2-es}
  v':\mathrm{P}^1_{C/A}\to\mathrm{P}^1_{C/B}
\end{equation}
tal que el diagrama
\[
\xymatrix{
  \mathrm{P}^1_{C/A} \ar[r]^{v'} & \mathrm{P}^1_{C/B} \\
  C \ar[u]^{p_1} \ar[r]^{\mathrm{I}_C} & C \ar[u]_{p_1}
}
\]
sea conmutativo. Como $v_0$ aplica $\mathfrak{J}_{C/A}$ en
$\mathfrak{J}_{C/B}$ se obtiene, restringiendo $v'$ a
$\Omega^1_{C/A}$, una aplicación
\begin{equation}
\tag{20.5.3.3}
\label{0.20.5.3.3-es}
  v_{C/B/A}:\Omega^1_{C/A}\to\Omega^1_{C/B}
\end{equation}
que es un homomorfismo de $C$-módulos.

\oldpage[0\textsubscript{IV-1}]{130}
Además, como el diagrama
\[
\xymatrix{
  \mathrm{P}^1_{C/A} \ar[r]^{v'} & \mathrm{P}^1_{C/B} \\
  C \ar[u]^{p_2} \ar[r]^{\mathrm{I}_C} & C \ar[u]_{p_2}
}
\tag{20.5.3.4}
\label{0.20.5.3.4-es}
\]
es también conmutativo, se deduce que el diagrama
\[
\xymatrix{
  \Omega^1_{C/A} \ar[r]^{v_{C/B/A}} & \Omega^1_{C/B} \\
  C \ar[u]^{d_{C/A}} \ar[r]^{\mathrm{I}_C} & C \ar[u]_{d_{C/B}}
}
\tag{20.5.3.5}
\label{0.20.5.3.5-es}
\]
es conmutativo.

Por último, si $s:A'\to A$ es un segundo homomorfismo de anillos, se tiene la
propiedad de transitividad
\begin{equation}
\tag{20.5.3.6}
\label{0.20.5.3.6-es}
  (v\circ s)_{C/B/A'}=v_{C/B/A}\circ s_{C/A/A'}.
\end{equation}
\end{env}

\begin{env}[20.5.4]
\label{0.20.5.4-es}
Si ahora se tiene un diagrama conmutativo de homomorfismos de anillos
\[
\xymatrix{
  B \ar[r]^v & B' \\
  A \ar[u]^f \ar[r]^u & A' \ar[u]_g
}
\]
se deduce de \sref{0.20.5.2.4-es} y \sref{0.20.5.3.3-es}, por composición,
un homomorfismo de $B'$-módulos
\begin{equation}
\tag{20.5.4.1}
\label{0.20.5.4.1-es}
  \Omega^1_{B/A}\otimes_B B'
  \xrightarrow{v_{B'/B/A}}\Omega^1_{B'/A}
  \xrightarrow{u_{B'/A'/A}}\Omega^1_{B'/A'}
\end{equation}
tal que el diagrama de $A'$-homomorfismos
\[
\xymatrix{
  \Omega^1_{B/A}\otimes_B B' \ar[r] & \Omega^1_{B'/A'} \\
  B' \ar[u]^{d_{B/A}\otimes1} \ar[r]^{\mathrm{I}_{B'}} &
  B' \ar[u]_{d_{B'/A'}}
}
\tag{20.5.4.2}
\label{0.20.5.4.2-es}
\]
sea conmutativo.

El homomorfismo \sref{0.20.5.4.1-es} corresponde además a un
dihomomorfismo de $B$-módulos
\begin{equation}
\tag{20.5.4.3}
\label{0.20.5.4.3-es}
  \Omega^1_{B/A}\to\Omega^1_{B'/A'}.
\end{equation}
\end{env}

\begin{proposition}[20.5.5]
\label{0.20.5.5-es}
Si $A'$, $B$ son dos $A$-álgebras y $B'=B\otimes_A A'$,
el homomorfismo canónico \sref{0.20.5.4.1-es}
\begin{equation}
\tag{20.5.5.1}
\label{0.20.5.5.1-es}
  \Omega^1_{B/A}\otimes_B B'\to\Omega^1_{B'/A'}
\end{equation}
es biyectivo.
\end{proposition}

\begin{proof}
\oldpage[0\textsubscript{IV-1}]{131}
El primer miembro de \sref{0.20.5.5.1-es} no es entonces otro que
$\Omega^1_{B/A}\otimes_A A'$. Se puede escribir
$B'\otimes_{A'}B'=(B\otimes_A B)\otimes_A A'$ salvo un isomorfismo canónico
y $p_{B'/A'}$ se identifica entonces con
$p_{B/A}\otimes\mathrm{I}_{A'}$; por consiguiente (como $p_{B/A}$ es sobreyectivo)
$\mathfrak{J}_{B'/A'}=\operatorname{Im}(\mathfrak{J}_{B/A}\otimes_A A')$
y
$\mathfrak{J}^2_{B'/A'}
=\operatorname{Im}(\mathfrak{J}^2_{B/A}\otimes_A A')$, de donde
$\mathrm{P}^1_{B'/A'}=\mathrm{P}^1_{B/A}\otimes_A A'$ salvo un isomorfismo
canónico, que hace corresponder entre sí los ideales de aumentación;
como el $A$-módulo $\mathrm{P}^1_{B/A}$ se identifica canónicamente con la
suma directa de $B$ y de $\Omega^1_{B/A}$, se tiene
\begin{equation}
\tag{20.5.5.2}
\label{0.20.5.5.2-es}
  \Omega^1_{B/A}\otimes_A A'=\Omega^1_{B'/A'}
\end{equation}
por el mismo isomorfismo, y se verifica inmediatamente que el compuesto de este
isomorfismo y del isomorfismo canónico
$\Omega^1_{B/A}\otimes_B B'\xrightarrow{\sim}
\Omega^1_{B/A}\otimes_A A'$ no es otro que \sref{0.20.5.5.1-es}.
\end{proof}

\begin{env}[20.5.6]
\label{0.20.5.6-es}
Los homomorfismos canónicos \sref{0.20.5.2.4-es} y
\sref{0.20.5.3.3-es} dan, por funtorialidad, para todo $C$-módulo $L$,
homomorfismos canónicos
\begin{align}
\tag{20.5.6.1}
\label{0.20.5.6.1-es}
  &\operatorname{Hom}_C(\Omega^1_{C/A},L)
  \to\operatorname{Hom}_C(\Omega^1_{B/A}\otimes_B C,L)
  =\operatorname{Hom}_B(\Omega^1_{B/A},L) \\
\tag{20.5.6.2}
\label{0.20.5.6.2-es}
  &\operatorname{Hom}_C(\Omega^1_{C/B},L)
  \to\operatorname{Hom}_C(\Omega^1_{C/A},L).
\end{align}
Teniendo en cuenta \sref{0.20.4.8.2-es} y los diagramas conmutativos
\sref{0.20.5.2.6-es} y \sref{0.20.5.3.5-es}, estos homomorfismos no son
otros (salvo una identificación canónica) que los homomorfismos
\sref{0.20.2.1.2-es} y \sref{0.20.2.1.3-es} respectivamente.
\end{env}

\begin{theorem}[20.5.7]
\label{0.20.5.7-es}
Sean $u:A\to B$, $v:B\to C$ dos homomorfismos de anillos.
\begin{enumerate}
  \item[\rm(i)] La sucesión de $C$-módulos
\begin{equation}
\tag{20.5.7.1}
\label{0.20.5.7.1-es}
  \Omega^1_{B/A}\otimes_B C
  \xrightarrow{v_{C/B/A}}\Omega^1_{C/A}
  \xrightarrow{u_{C/B/A}}\Omega^1_{C/B}\to0
\end{equation}
es exacta.
  \item[\rm(ii)] Para que $v_{C/B/A}$ sea invertible por la izquierda, es
  necesario y suficiente que $C$ sea una $B$-álgebra formalmente lisa relativamente a
  $A$ (para las topologías discretas (cf. \sref{0.19.9.1-es})); en
  particular, basta para ello que $C$ sea una $B$-álgebra formalmente
  lisa (para las topologías discretas).
\end{enumerate}
\end{theorem}

\begin{proof}
(i) La exactitud de la sucesión \sref{0.20.2.4.2-es} muestra en primer lugar,
teniendo en cuenta \sref{0.20.5.6-es}, que la sucesión
\[
  0\to\operatorname{Hom}_C(\Omega^1_{C/B},L)
  \to\operatorname{Hom}_C(\Omega^1_{C/A},L)
  \to\operatorname{Hom}_C(\Omega^1_{B/A}\otimes_B C,L)
\]
es exacta para todo $C$-módulo $L$. Se sabe que esto implica la exactitud
de la sucesión \sref{0.20.5.7.1-es} (Bourbaki, \emph{Alg.}, cap.~II,
3\textsuperscript{a} ed., §~2, n\textsuperscript{o}~1, teo.~1).

(ii) En virtud de la exactitud de \sref{0.20.5.7.1-es}, decir que
$v_{C/B/A}$ es invertible por la izquierda significa que la sucesión
\begin{equation}
\tag{20.5.7.2}
\label{0.20.5.7.2-es}
  0\to\Omega^1_{B/A}\otimes_B C
  \xrightarrow{v_{C/B/A}}\Omega^1_{C/A}
  \xrightarrow{u_{C/B/A}}\Omega^1_{C/B}\to0
\end{equation}
es \emph{exacta y escindida}; se sabe (Bourbaki, \emph{loc. cit.},
n\textsuperscript{o}~1, prop.~1) que esto equivale a decir que para todo
$C$-módulo $L$, la sucesión
\[
  0\to\operatorname{Hom}_C(\Omega^1_{C/B},L)
  \to\operatorname{Hom}_C(\Omega^1_{C/A},L)
  \to\operatorname{Hom}_C(\Omega^1_{B/A}\otimes_B C,L)\to0
\]
es exacta; teniendo en cuenta \sref{0.20.5.6-es} y
\sref{0.20.2.4.2-es}, esta condición equivale a
\[
  \operatorname{Exalcom}_{B/A}(C,L)=0
\]
para todo $C$-módulo $L$, y la conclusión resulta por tanto de
\sref{0.19.9.8.1-es}.
\end{proof}

\oldpage[0\textsubscript{IV-1}]{132}
Observemos además que si se tiene un diagrama conmutativo de homomorfismos
de anillos
\[
\xymatrix{
  A' \ar[r] & B' \ar[r] & C' \\
  A \ar[u] \ar[r] & B \ar[u] \ar[r] & C \ar[u]
}
\]
se deduce un diagrama conmutativo
\[
\xymatrix{
  \Omega^1_{B/A}\otimes_B C \ar[r] \ar[d] &
  \Omega^1_{C/A} \ar[r] \ar[d] &
  \Omega^1_{C/B} \ar[r] \ar[d] & 0 \\
  \Omega^1_{B'/A'}\otimes_{B'} C' \ar[r] &
  \Omega^1_{C'/A'} \ar[r] &
  \Omega^1_{C'/B'} \ar[r] & 0
}
\tag{20.5.7.3}
\label{0.20.5.7.3-es}
\]
donde las flechas verticales provienen de los dihomomorfismos
\sref{0.20.5.4.3-es}.

\begin{corollary}[20.5.8]
\label{0.20.5.8-es}
Supongamos que el homomorfismo $v:B\to C$ haga de $C$ una $B$-álgebra
formalmente étale (para las topologías discretas
\sref{0.19.10.2-es}); entonces el homomorfismo \sref{0.20.5.3.3-es}
\[
  v_{C/B/A}:\Omega^1_{B/A}\otimes_B C\to\Omega^1_{C/A}
\]
es biyectivo.
\end{corollary}

\begin{proof}
En efecto, si $C$ es una $B$-álgebra formalmente no ramificada para las
topologías discretas, se deduce de \sref{0.19.10.4-es},
\sref{0.20.4.8-es} y \sref{0.20.1.1-es} que se tiene
$\operatorname{Hom}_C(\Omega^1_{C/B},L)=0$ para todo $C$-módulo $L$, por tanto
$\Omega^1_{C/B}=0$ (cf. (20.7.4)); por otra parte, si $C$ es una
$B$-álgebra formalmente lisa para las topologías discretas, la sucesión
\sref{0.20.5.7.2-es} es exacta; de donde resulta el corolario.
\end{proof}

\begin{corollary}[20.5.9]
\label{0.20.5.9-es}
Sean $A$ un anillo, $B$ una $A$-álgebra, $S$ un subconjunto multiplicativo de
$B$; entonces el homomorfismo canónico
\begin{equation}
\tag{20.5.9.1}
\label{0.20.5.9.1-es}
  S^{-1}\Omega^1_{B/A}\to\Omega^1_{S^{-1}B/A}
\end{equation}
es biyectivo.
\end{corollary}

\begin{proof}
Basta aplicar \sref{0.20.5.8-es} a $C=S^{-1}B$, que es una
$B$-álgebra formalmente étale para las topologías discretas
\sref{0.19.10.3-es}, (ii).
\end{proof}

Teniendo en cuenta \sref{0.20.5.5-es}, se puede por tanto escribir
\begin{equation}
\tag{20.5.9.2}
\label{0.20.5.9.2-es}
  \Omega^1_{S^{-1}B/S^{-1}A}
  =S^{-1}\Omega^1_{B/A}
  =\Omega^1_{S^{-1}B/A},
\end{equation}
salvo isomorfismos canónicos.

\begin{corollary}[20.5.10]
\label{0.20.5.10-es}
Si $k$ es un cuerpo y $K=k(X_\alpha)_{\alpha\in I}$ una extensión puramente trascendente de
$k$, los $dX_\alpha$ forman una base del $K$-espacio vectorial
$\Omega^1_{K/k}$.
\end{corollary}

\begin{proof}
Como $K$ es el cuerpo de fracciones del anillo de polinomios
$k[X_\alpha]_{\alpha\in I}$, esto resulta de \sref{0.20.4.13-es}, (i) y de
\sref{0.20.5.9-es}.
\end{proof}

\oldpage[0\textsubscript{IV-1}]{133}
\begin{env}[20.5.11]
\label{0.20.5.11-es}
Sean $A$ un anillo, $B$ una $A$-álgebra, $\mathfrak{R}$ un ideal de $B$,
$C$ la $A$-álgebra cociente $B/\mathfrak{R}$, y consideremos
el homomorfismo compuesto de $A$-módulos
\begin{equation}
\tag{20.5.11.1}
\label{0.20.5.11.1-es}
  \delta:\mathfrak{R}\longrightarrow B\xrightarrow{d_{B/A}}\Omega^1_{B/A}
\end{equation}
donde la primera flecha es la inyección canónica; como
$d(xy)=xdy+ydx$, se ve que se tiene
$\delta(\mathfrak{R}^2)\subset\mathfrak{R}.\Omega^1_{B/A}$, de donde, por
paso a los cocientes, un homomorfismo de $A$-módulos
\begin{equation}
\tag{20.5.11.2}
\label{0.20.5.11.2-es}
  \delta_{C/B/A}:\mathfrak{R}/\mathfrak{R}^2\longrightarrow
  \Omega^1_{B/A}\otimes_B C=
  \Omega^1_{B/A}/(\mathfrak{R}.\Omega^1_{B/A}).
\end{equation}
Pero de hecho, $\delta_{C/B/A}$ es un homomorfismo de $C$-módulos, pues
para $x\in B$ e $y\in\mathfrak{R}$, se tiene
$ydx\in\mathfrak{R}.\Omega^1_{B/A}$, por tanto
$d(xy)\equiv xdy\pmod{\mathfrak{R}.\Omega^1_{B/A}}$, lo que prueba
primero que (20.5.11.2) es un homomorfismo de $B$-módulos, y como
$\mathfrak{R}$ anula los dos miembros, esto establece nuestra afirmación.

Si $B'$ es una segunda $A$-álgebra, $u:B\to B'$ un $A$-homomorfismo,
$\mathfrak{R}'$ un ideal de $B'$ tal que
$u(\mathfrak{R})\subset\mathfrak{R}'$ y $C'=B'/\mathfrak{R}'$ el álgebra
cociente, se tiene un diagrama conmutativo
\begin{equation}
\tag{20.5.11.3}
\label{0.20.5.11.3-es}
\xymatrix{
  \mathfrak{R}/\mathfrak{R}^2 \ar[r]^-{\delta_{C/B/A}} \ar[d] &
  \Omega^1_{B/A}\otimes_B C \ar[d] \\
  \mathfrak{R}'/\mathfrak{R}'^2 \ar[r]_-{\delta_{C'/B'/A}} &
  \Omega^1_{B'/A}\otimes_{B'} C'
}
\end{equation}
donde las flechas verticales provienen de $u$ (20.5.2.4).
\end{env}

\begin{theorem}[20.5.12]
\label{0.20.5.12-es}
Sean $A$ un anillo, $B$ una $A$-álgebra, $C$ la $A$-álgebra cociente
$B/\mathfrak{R}$, $u:B\to C$ el homomorfismo canónico.
\begin{enumerate}
  \item[\rm(i)] Se tiene una sucesión exacta de $C$-módulos
\begin{equation}
\tag{20.5.12.1}
\label{0.20.5.12.1-es}
  \mathfrak{R}/\mathfrak{R}^2
  \xrightarrow{\delta_{C/B/A}}\Omega^1_{B/A}\otimes_B C
  \xrightarrow{u_{C/B/A}}\Omega^1_{C/A}\longrightarrow 0,
\end{equation}
donde $u_{C/B/A}$ y $\delta_{C/B/A}$ están definidos por (20.5.2.4) y
(20.5.11.2) respectivamente.
  \item[\rm(ii)] Si se pone $E=B/\mathfrak{R}^2$, el homomorfismo
canónico (20.5.2.4)
\[
  \Omega^1_{B/A}\otimes_B C\longrightarrow
  \Omega^1_{E/A}\otimes_E C
\]
es biyectivo.
  \item[\rm(iii)] Las tres condiciones siguientes son equivalentes:
  \begin{enumerate}
    \item[\rm a)] $\delta_{C/B/A}$ es invertible por la izquierda.
    \item[\rm b)] Toda $A$-extensión de $C$ por un $C$-módulo $L$, cuya
imagen inversa por $u:B\to C$ es $A$-trivial, es ella misma
$A$-trivial.
    \item[\rm c)] La $A$-álgebra $E=B/\mathfrak{R}^2$ es una extensión
$A$-trivial de $C$ por $\mathfrak{R}/\mathfrak{R}^2$.
  \end{enumerate}
  \item[\rm(iv)] Hay una correspondencia biunívoca canónica entre los
inversos por la izquierda de $\delta_{C/B/A}$ y los inversos por la derecha del
homomorfismo canónico $E\to C$.
\end{enumerate}
\end{theorem}

\begin{proof}
\oldpage[0\textsubscript{IV-1}]{134}
(i) Como $u$ es sobreyectivo, se tiene $\operatorname{D\acute{e}r}_B(C,L)=0$ para todo $C$-módulo
$L$ en virtud de (20.1.1). La sucesión exacta (20.2.3.1) se convierte por tanto en
\begin{equation}
\tag{20.5.12.2}
\label{0.20.5.12.2-es}
  0\longrightarrow\operatorname{D\acute{e}r}_A(C,L)\xrightarrow{v^0}\operatorname{D\acute{e}r}_A(B,L)
  \xrightarrow{\partial}\operatorname{Exalcom}_B(C,L)
  \xrightarrow{v^1}\operatorname{Exalcom}_A(C,L)
  \xrightarrow{u^1}\operatorname{Exalcom}_A(B,L),
\end{equation}
donde $v$ es el homomorfismo $A\to B$. Recordemos, por otra parte, (18.3.8) que
$\operatorname{Exalcom}_B(C,L)$ se identifica canónicamente con
$\operatorname{Hom}_C(\mathfrak{R}/\mathfrak{R}^2,L)$; se deduce, pues, de (20.5.12.2)
y (20.4.8) la sucesión exacta
\begin{equation}
\tag{20.5.12.3}
\label{0.20.5.12.3-es}
\begin{split}
  0\longrightarrow\operatorname{Hom}_C(\Omega^1_{C/A},L)
  \longrightarrow\operatorname{Hom}_C(\Omega^1_{B/A}\otimes_B C,L)
  \xrightarrow{\varphi}\operatorname{Hom}_C(\mathfrak{R}/\mathfrak{R}^2,L)\\
  \xrightarrow{\psi}\operatorname{Ker}(u^1)\longrightarrow 0
\end{split}
\end{equation}
con $\varphi=\partial\circ\eta^{-1}$ y $\psi=v^1\circ\eta$. Volviendo
a las definiciones de $\partial$ (20.2.2) y de $\eta$ (18.3.8), se comprueba
inmediatamente que $\varphi$ es precisamente el homomorfismo
$\operatorname{Hom}(\delta_{C/B/A},1_L)$. La existencia de la sucesión exacta formada por los
4 primeros términos de (20.5.12.3) muestra, pues, que la sucesión (20.5.12.1)
es exacta (Bourbaki, \emph{Álg.}, cap. II, 3\textsuperscript{a} ed.,
\S~2, n\textsuperscript{o}~1, teorema~1).

\medskip
(ii) Apliquemos a $B$ y al ideal $\mathfrak{R}^2$ la sucesión exacta
(20.5.12.1), lo que da
\begin{equation}
\tag{20.5.12.4}
\label{0.20.5.12.4-es}
  \mathfrak{R}^2/\mathfrak{R}^4\longrightarrow
  \Omega^1_{B/A}\otimes_B E\longrightarrow\Omega^1_{E/A}\longrightarrow 0,
\end{equation}
de donde, tensorizando por $C$ (considerado como $E$-álgebra), se obtiene la sucesión
exacta
\[
  \mathfrak{R}^2/\mathfrak{R}^4\otimes_E C\longrightarrow
  \Omega^1_{B/A}\otimes_B C\longrightarrow
  \Omega^1_{E/A}\otimes_E C\longrightarrow 0.
\]
Ahora bien, si $x,y$ son dos elementos de $\mathfrak{R}$, y $\zeta$ la clase
de $xy$ mód. $\mathfrak{R}^4$, la imagen
$\delta'(\zeta\otimes1)$ es, por definición,
$d_{B/A}(xy)\otimes1=(xd_{B/A}(y)+yd_{B/A}(x))\otimes1$; pero, como las
imágenes de $x$ y de $y$ en $C$ son nulas, se tiene también
$d_{B/A}(xy)\otimes1=0$ en $\Omega^1_{B/A}\otimes_B C$, lo que demuestra
la aserción.

\medskip
(iii) Decir que $\delta_{C/B/A}$ es invertible por la izquierda equivale, teniendo
en cuenta la exactitud de (20.5.12.1), a decir que la sucesión
\begin{equation}
\tag{20.5.12.5}
\label{0.20.5.12.5-es}
  0\longrightarrow\mathfrak{R}/\mathfrak{R}^2
  \xrightarrow{\delta_{C/B/A}}\Omega^1_{B/A}\otimes_B C
  \xrightarrow{u_{C/B/A}}\Omega^1_{C/A}\longrightarrow0
\end{equation}
es \emph{exacta y escindida}, y equivale también (Bourbaki,
\emph{loc. cit.}) a decir que $\operatorname{Ker}(u^1)=0$ en la sucesión exacta
(20.5.12.3) para todo $L$, lo que muestra la equivalencia de las condiciones
a) y b) (cf. (18.3.6.2)).

El hecho de que b) implique c) proviene de que la imagen inversa por
$u:B\to C$ de la $A$-extensión $E$ de $C$ por
$\mathfrak{R}/\mathfrak{R}^2$ es $A$-trivial, pues $u$ es la composición de los
homomorfismos canónicos $B\to E\to C$ (18.1.6). Recíprocamente, c)
implica b), pues toda $B$-extensión de $C$ por $L$ es $B$-equivalente
a $E\oplus_{\mathfrak{R}/\mathfrak{R}^2}L$ (18.3.8), es decir, su
clase es la imagen de la clase de $E$ (considerada como $B$-extensión)
por el homomorfismo

\oldpage[0\textsubscript{IV-1}]{135}
$w_*:\operatorname{Exan}_B(C,\mathfrak{R}/\mathfrak{R}^2)
\to\operatorname{Exan}_B(C,L)$ correspondiente a un $B$-homomorfismo
$w:\mathfrak{R}/\mathfrak{R}^2\to L$. El hecho de que c) implique b) se deduce
entonces de la conmutatividad del diagrama (18.3.6.5)
\[
\xymatrix{
  \operatorname{Exan}_B(C,\mathfrak{R}/\mathfrak{R}^2)
    \ar[r]^{w_*} \ar[d]_{u^*} &
  \operatorname{Exan}_B(C,L) \ar[d]^{u^*} \\
  \operatorname{Exan}_A(C,\mathfrak{R}/\mathfrak{R}^2)
    \ar[r]_{w_*} & \operatorname{Exan}_A(C,L).
}
\]

\medskip
(iv) Se ha visto (20.1.7) que los inversos por la derecha de $E\to C$
corresponden canónica y biyectivamente al conjunto de los elementos
$D\in\operatorname{D\acute{e}r}_A(E,\mathfrak{R}/\mathfrak{R}^2)$ tales que $D(x)=x$ en
$\mathfrak{R}/\mathfrak{R}^2$, y por tanto también, por (20.4.8), al conjunto de los
$E$-homomorfismos
$h:\Omega^1_{E/A}\to\mathfrak{R}/\mathfrak{R}^2$ tales que la composición
\[
  \mathfrak{R}/\mathfrak{R}^2\xrightarrow{d_{E/A}}
  \Omega^1_{E/A}\xrightarrow{h}\mathfrak{R}/\mathfrak{R}^2
\]
sea la identidad. Tensorizando con $C$, se deduce de ello (puesto que
$\mathfrak{R}/\mathfrak{R}^2$ es un $C$-módulo) que la composición
\[
  \mathfrak{R}/\mathfrak{R}^2\xrightarrow{\delta_{C/E/A}}
  \Omega^1_{E/A}\otimes_E C\xrightarrow{h\otimes1}
  \mathfrak{R}/\mathfrak{R}^2
\]
es la identidad; ahora bien, puesto que $\mathfrak{R}/\mathfrak{R}^2$ es un
$C$-módulo, $h\mapsto h\otimes1$ es un isomorfismo del conjunto
$\operatorname{Hom}_E(\Omega^1_{E/A},\mathfrak{R}/\mathfrak{R}^2)$ sobre
$\operatorname{Hom}_C(\Omega^1_{E/A}\otimes_E C,\mathfrak{R}/\mathfrak{R}^2)$; por
otra parte, (ii) demuestra que se pueden identificar canónicamente
$\Omega^1_{E/A}\otimes_E C$ y $\Omega^1_{B/A}\otimes_B C$,
identificándose entonces $\delta_{C/E/A}$ con $\delta_{C/B/A}$. C.Q.D.
\end{proof}

\begin{example}[20.5.13]
\label{0.20.5.13-es}
Sean $B=A[X_\alpha]_{\alpha\in I}$ un álgebra de polinomios sobre $A$,
$\mathfrak{R}$ un ideal de $B$, $(P_\lambda)$ un sistema de generadores
de $\mathfrak{R}$ y $C=B/\mathfrak{R}$; se sabe que
$\Omega^1_{B/A}$ es un $B$-módulo libre y que los $dX_\alpha$ forman una
base (20.4.13, (i)); por tanto, los $dX_\alpha$ forman también una base del
$C$-módulo libre $\Omega^1_{B/A}\otimes_B C$. Por otra parte, se deduce
inmediatamente de la definición que la imagen de
$\mathfrak{R}/\mathfrak{R}^2$ por $\delta_{C/B/A}$ es el sub-$C$-módulo
generado por los
\[
  dP_\lambda=\sum_\alpha
  \frac{\partial P_\lambda}{\partial X_\alpha}\,dX_\alpha.
\]
Se concluye que $\Omega^1_{C/A}$ es isomorfo al cociente del
$C$-módulo libre que tiene los $dX_\alpha$ por base, por el sub-$C$-módulo
generado por los $dP_\lambda$, lo que da una descripción de un módulo
de diferenciales de un álgebra \emph{cualquiera}, pues toda $A$-álgebra
$C$ puede obtenerse de la manera precedente.
\end{example}

\begin{corollary}[20.5.14]
\label{0.20.5.14-es}
Si $C$ es un $A$-álgebra formalmente lisa (para las topologías
discretas), la sucesión
\begin{equation}
\tag{20.5.14.1}
\label{0.20.5.14.1-es}
  0\longrightarrow\mathfrak{R}/\mathfrak{R}^2
  \xrightarrow{\delta_{C/B/A}}\Omega^1_{B/A}\otimes_B C
  \xrightarrow{u_{C/B/A}}\Omega^1_{C/A}\longrightarrow0
\end{equation}
es exacta y escindida.
\end{corollary}

\begin{proof}
En efecto, toda $A$-extensión de $C$ por un $C$-módulo es entonces trivial
(19.4.4.1).
\end{proof}

\oldpage[0\textsubscript{IV-1}]{136}
\begin{remark}[20.5.15]
\label{0.20.5.15-es}
Sea $u:A\to B$ un homomorfismo \emph{sobreyectivo} de anillos; entonces, para
todo homomorfismo de anillos $v:B\to C$, el homomorfismo canónico
\begin{equation}
\tag{20.5.15.1}
\label{0.20.5.15.1-es}
  u_{C/B/A}:\Omega^1_{C/A}\longrightarrow\Omega^1_{C/B}
\end{equation}
es \emph{biyectivo}; esto se deduce, en efecto, de la sucesión exacta (20.5.7.1),
puesto que $\Omega^1_{B/A}=0$ (20.4.12).
\end{remark}

\subsection{Módulos de imperfección y homomorfismos característicos}
\label{subsection:0.20.6-es}
\label{0.20.6-es}

\begin{definition}[20.6.1]
\label{0.20.6.1-es}
Dados dos homomorfismos de anillos $u:A\to B$, $v:B\to C$, se
denomina \emph{módulo de imperfección del $B$-álgebra $C$ relativo a
$A$}, y se denota $\Upsilon_{C/B/A}$, al $C$-módulo núcleo del
homomorfismo
$v_{C/B/A}:\Omega^1_{B/A}\otimes_B C\to\Omega^1_{C/A}$.
Se tiene, pues, por definición (cf. (20.5.7)) la sucesión exacta
\begin{equation}
\tag{20.6.1.1}
\label{0.20.6.1.1-es}
  0\longrightarrow\Upsilon_{C/B/A}\longrightarrow
  \Omega^1_{B/A}\otimes_B C\xrightarrow{v_{C/B/A}}
  \Omega^1_{C/A}\xrightarrow{u_{C/B/A}}\Omega^1_{C/B}
  \longrightarrow0.
\end{equation}
\end{definition}

Cuando $A=\mathbb{Z}$ (siendo entonces los módulos $\Omega^1_{B/A}$ y $\Omega^1_{C/A}$
los módulos de diferenciales «absolutas» $\Omega^1_B$ y
$\Omega^1_C$), se escribe $\Upsilon_{C/B}$ en lugar de
$\Upsilon_{C/B/\mathbb{Z}}$; cuando $B$ y $C$ son álgebras sobre un
cuerpo primo $P$, se tiene $\Upsilon_{C/B/P}=\Upsilon_{C/B}$.

Sean $R$, $S$ partes multiplicativas de $B$ y $C$, respectivamente,
tales que la imagen de $R$ esté contenida en $S$. Se deduce entonces de la
sucesión exacta (20.6.1.1) y de (20.5.9) que se tiene
\begin{equation}
\tag{20.6.1.2}
\label{0.20.6.1.2-es}
  \Upsilon_{S^{-1}C/R^{-1}B/A}=S^{-1}\Upsilon_{C/B/A}.
\end{equation}

\begin{proposition}[20.6.2]
\label{0.20.6.2-es}
Si $C$ es un $B$-álgebra formalmente lisa relativamente a $A$ (y, en
particular, si $C$ es un $B$-álgebra formalmente lisa), se tiene
$\Upsilon_{C/B/A}=0$.
\end{proposition}

\begin{proof}
Esto se deduce de (20.5.7, (ii)).
\end{proof}

\begin{proposition}[20.6.3]
\label{0.20.6.3-es}
Sean $k$ un cuerpo y $K$ una extensión de $k$. Para que $K$ sea una
extensión separable de $k$, es necesario y suficiente que $\Upsilon_{K/k}=0$
(es decir, que el homomorfismo canónico
$\Omega^1_k\otimes_k K\to\Omega^1_K$ sea inyectivo o también (20.4.8), que
toda derivación de $k$ en $K$ se prolongue a una derivación de $K$
en sí mismo).
\end{proposition}

\begin{proof}
En efecto, equivale decir que $K$ es separable sobre $k$ o que es un
$k$-álgebra formalmente lisa (19.6.1); por otra parte, si $P$ es el
cuerpo primo de $k$, $K$ es separable sobre $P$, y por tanto un $P$-álgebra
formalmente lisa, y equivale, pues, decir que $K$ es un
$k$-álgebra formalmente lisa o un $k$-álgebra formalmente lisa
relativamente a $P$. Finalmente, decir que $K$ es un $k$-álgebra formalmente
lisa relativamente a $P$ equivale, según (20.5.7, (ii)), a decir que
el homomorfismo
$v_{K/k/P}:\Omega^1_k\otimes_k K\to\Omega^1_K$ es invertible por la izquierda;
pero, como $K$ es un cuerpo, esta última condición equivale a decir
que el núcleo de $v_{K/k/P}$, es decir, $\Upsilon_{K/k}$, es nulo.
\end{proof}

\begin{env}[20.6.4]
\label{0.20.6.4-es}
Consideremos un diagrama conmutativo
\begin{equation}
\tag{20.6.4.1}
\label{0.20.6.4.1-es}
\xymatrix{
  A' \ar[r]^{u'} & B' \ar[r]^{v'} & C' \\
  A \ar[u]^f \ar[r]_u & B \ar[u]^g \ar[r]_v & C \ar[u]^h
}
\end{equation}

\oldpage[0\textsubscript{IV-1}]{137}
de homomorfismos de anillos conmutativos. La conmutatividad del diagrama
correspondiente (20.5.7.3) implica la existencia de un único
$C$-homomorfismo
\begin{equation}
\tag{20.6.4.2}
\label{0.20.6.4.2-es}
  \Upsilon_{C/B/A}\longrightarrow\Upsilon_{C'/B'/A'}
\end{equation}
deducido canónicamente de (20.6.4.1) y que hace conmutativo el diagrama
\begin{equation}
\tag{20.6.4.3}
\label{0.20.6.4.3-es}
\xymatrix{
  0 \ar[r] & \Upsilon_{C/B/A} \ar[r] \ar[d] &
  \Omega^1_{B/A}\otimes_B C \ar[r]^{v_{C/B/A}} \ar[d] &
  \Omega^1_{C/A} \ar[r]^{u_{C/B/A}} \ar[d] &
  \Omega^1_{C/B} \ar[r] \ar[d] & 0 \\
  0 \ar[r] & \Upsilon_{C'/B'/A'} \ar[r] &
  \Omega^1_{B'/A'}\otimes_{B'} C' \ar[r]_{v_{C'/B'/A'}} &
  \Omega^1_{C'/A'} \ar[r]_{u_{C'/B'/A'}} &
  \Omega^1_{C'/B'} \ar[r] & 0
}
\end{equation}

El dato del homomorfismo (20.6.4.2) equivale además al de un
$C'$-homomorfismo
\begin{equation}
\tag{20.6.4.4}
\label{0.20.6.4.4-es}
  \Upsilon_{C/B/A}\otimes_C C'\longrightarrow\Upsilon_{C'/B'/A'}
\end{equation}
que, compuesto con el homomorfismo canónico
$\Upsilon_{C/B/A}\to\Upsilon_{C/B/A}\otimes_C C'$, vuelve a dar (20.6.4.2).
Está claro que (20.6.4.2) goza de una evidente propiedad de
transitividad, lo que permite decir que $\Upsilon_{C/B/A}$ es un
\emph{funtor} de la terna $(A,B,C)$.
\end{env}

\begin{env}[20.6.5]
\label{0.20.6.5-es}
Nos resultará conveniente para lo que sigue, bajo las condiciones de (20.6.1),
introducir un \emph{complejo} (de cadenas) \emph{de $C$-módulos}
$\mathrm{K}_\bullet(C/B/A)$ cuyos términos son nulos salvo en grados 0
y 1, donde se toma
\begin{equation}
\tag{20.6.5.1}
\label{0.20.6.5.1-es}
\left\{
\begin{aligned}
  \mathrm{K}_0(C/B/A)&=\Omega^1_{C/A},\\
  \mathrm{K}_1(C/B/A)&=\Omega^1_{B/A}\otimes_B C,
\end{aligned}
\right.
\end{equation}
siendo el operador diferencial $\mathrm{K}_1\to\mathrm{K}_0$
$v_{C/B/A}$. Esto permite escribir (salvo isomorfismos canónicos)
$\Omega^1_{C/B}$ y $\Upsilon_{C/B/A}$ como los \emph{módulos
de homología} de este complejo:
\begin{equation}
\tag{20.6.5.2}
\label{0.20.6.5.2-es}
  \mathrm{H}_0(\mathrm{K}_\bullet(C/B/A))=\Omega^1_{C/B},\qquad
  \mathrm{H}_1(\mathrm{K}_\bullet(C/B/A))=\Upsilon_{C/B/A}.
\end{equation}
\end{env}

Del mismo modo:
\begin{proposition}[20.6.6]
\label{0.20.6.6-es}
Bajo las hipótesis de (20.6.1), para todo $C$-módulo $L$, se tienen
$C$-isomorfismos canónicos
\begin{align}
\tag{20.6.6.1}
\label{0.20.6.6.1-es}
  \mathrm{H}^0(\mathrm{K}_\bullet(C/B/A),L)&\simeq\operatorname{D\acute{e}r}_B(C,L),\\
\tag{20.6.6.2}
\label{0.20.6.6.2-es}
  \mathrm{H}^1(\mathrm{K}_\bullet(C/B/A),L)&\simeq
  \operatorname{Exalcom}_{B/A}(C,L).
\end{align}
\end{proposition}

\begin{proof}
En efecto, el complejo de cocadenas
$\operatorname{Hom}_C(\mathrm{K}_\bullet(C/B/A),L)$ no es otro, en virtud de
(20.4.8)
y (20.5.6), que el complejo
\[
  \cdots\longrightarrow0\longrightarrow\operatorname{D\acute{e}r}_A(C,L)
  \longrightarrow\operatorname{D\acute{e}r}_A(B,L)\longrightarrow0\longrightarrow\cdots
\]
\oldpage[0\textsubscript{IV-1}]{138}
donde el operador diferencial es $v^0$ (con las notaciones de (20.2.1)).
La proposición resulta entonces de la sucesión exacta (20.2.4.2) y de la
definición de los módulos de cohomología
\begin{equation}
\tag{20.6.6.3}
\label{0.20.6.6.3-es}
  \mathrm{H}^i(\mathrm{K}_\bullet,L)=
  \mathrm{H}^i(\operatorname{Hom}_C(\mathrm{K}_\bullet,L)).
\end{equation}
\end{proof}

Si se tiene un diagrama conmutativo de homomorfismos de anillos
\[
\xymatrix{
  A' \ar[r] & B' \ar[r] & C' \\
  A \ar[u] \ar[r] & B \ar[u] \ar[r] & C \ar[u]
}
\]
los dihomomorfismos (20.5.4.3) definen un dihomomorfismo de
complejos de módulos
\begin{equation}
\tag{20.6.6.4}
\label{0.20.6.6.4-es}
  \mathrm{K}_\bullet(C/B/A)\longrightarrow
  \mathrm{K}_\bullet(C'/B'/A')
\end{equation}
y los dihomomorfismos que de él se deducen para la homología o la
cohomología se identifican, por las fórmulas (20.6.5.2), (20.6.6.1) y
(20.6.6.2), con los dihomomorfismos ya definidos en (20.5.4.3),
(20.6.4.2), (20.2.1) y (18.3.7.2) (teniendo en cuenta, para el último, la
definición del operador $\partial$ en la sucesión exacta (20.2.4.2)).

\begin{env}[20.6.7]
\label{0.20.6.7-es}
Se sabe que, para un complejo $\mathrm{K}_\bullet$ de $C$-módulos y un
$C$-módulo $L$, se tienen homomorfismos canónicos
\[
  \alpha_i:\mathrm{H}^i(\mathrm{K}_\bullet,L)\longrightarrow
  \operatorname{Hom}_C(\mathrm{H}_i(\mathrm{K}_\bullet),L)
\]
(M, IV, 6). Aquí, el homomorfismo canónico
\[
  \alpha_1:\mathrm{H}^1(\mathrm{K}_\bullet(C/B/A),L)
  \longrightarrow
  \operatorname{Hom}_C(\mathrm{H}_1(\mathrm{K}_\bullet(C/B/A)),L)
\]
se define inmediatamente haciendo pasar al cociente por la imagen
de $\operatorname{Hom}_C(\mathrm{K}_0,L)$ el homomorfismo de restricción
\[
  \operatorname{Hom}_C(\mathrm{K}_1(C/B/A),L)\longrightarrow
  \operatorname{Hom}_C(\mathrm{H}_1(\mathrm{K}_\bullet(C/B/A)),L)
\]
puesto que $\mathrm{H}_1(\mathrm{K}_\bullet(C/B/A))$ no es otro que el
núcleo de $\mathrm{K}_1\to\mathrm{K}_0$; teniendo en cuenta (20.6.6.2) y
(20.6.5.2), se obtiene, pues, un $C$-homomorfismo canónico
\begin{equation}
\tag{20.6.7.1}
\label{0.20.6.7.1-es}
  \operatorname{Exalcom}_{B/A}(C,L)\longrightarrow
  \operatorname{Hom}_C(\Upsilon_{C/B/A},L)
\end{equation}
que se explicita de la manera siguiente: en virtud de (20.2.4.2), toda
$B$-extensión de $C$ por $L$ que es $A$-trivial procede del dato
de una $A$-derivación $D$ de $B$ en $L$, luego, por (20.4.8), de un
$C$-homomorfismo $f$ de $\Omega^1_{B/A}\otimes_B C$ en $L$; se hace
corresponder a la clase de esta extensión la \emph{restricción} de $f$
a $\Upsilon_{C/B/A}$, que solo depende, en efecto, de esta clase y
no de la elección de $D$.
\end{env}

\begin{definition}[20.6.8]
\label{0.20.6.8-es}
Sean $u:A\to B$, $v:B\to C$ dos homomorfismos de anillos. Para toda
$B$-extensión $E$ de $C$ por un $C$-módulo $L$, que sea $A$-trivial
(18.3.7), se llama \emph{homomorfismo característico de $E$}, y se
denota por $\chi_E$, al $C$-homomorfismo $\Upsilon_{C/B/A}\to L$, imagen de la
clase de $E$ por el homomorfismo canónico (20.6.7.1).
\end{definition}

\oldpage[0\textsubscript{IV-1}]{139}
Se puede definir el homomorfismo $\chi_E$ de otra manera:
\begin{proposition}[20.6.9]
\label{0.20.6.9-es}
Sea $E$ una $B$-extensión de $C$ por un $C$-módulo $L$, que sea
$A$-trivial (18.3.7); entonces el diagrama
\begin{equation}
\tag{20.6.9.1}
\label{0.20.6.9.1-es}
\xymatrix{
  0 \ar[r] & \Upsilon_{C/B/A} \ar[r] \ar[d]^{\chi_E} &
  \Omega^1_{B/A}\otimes_B C \ar[r]^{v_{C/B/A}}
    \ar[d]^{q_{E/B/A}\otimes1_C} &
  \Omega^1_{C/A} \ar[d]^{1} \\
  0 \ar[r] & L \ar[r]_{\delta_{C/E/A}} &
  \Omega^1_{E/A}\otimes_E C \ar[r]_{p_{C/E/A}} &
  \Omega^1_{C/A} \ar[r] & 0
}
\end{equation}
en el que $q:B\to E$ define la estructura de $B$-extensión de $E$ y
$p:E\to C$ es el homomorfismo de aumentación, es conmutativo; sus
filas son exactas.
\end{proposition}

\begin{proof}
La fila inferior del diagrama es la sucesión (20.5.12.5) relativa a los
dos homomorfismos $A\xrightarrow{q\circ u}E$ y
$p:E\to C=E/L$; como $p$ es sobreyectivo y $E$ es una extensión
$A$-trivial, esta sucesión es exacta y escindida en virtud de (20.5.12,
(iii)). La conmutatividad del cuadrado derecho de (20.6.9.1) resulta de la
relación $v=p\circ q$ (20.5.2.7); la imagen por
$q_{E/B/A}\otimes1_C$ del núcleo $\Upsilon_{C/B/A}$ de $v_{C/B/A}$ está
por tanto contenida en el núcleo $L$ de $p_{C/E/A}$. Por otra parte, sea
$h:C\to E$ un $A$-homomorfismo inverso por la derecha de $p$, y sea
$j:L\to E$ la inyección canónica, de modo que se tiene
$q(b)=h(v(b))+j(D(b))$ para $b\in B$, donde $D$ es la $A$-derivación de
$B$ en $L$ que define la $B$-extensión $E$; se puede escribir
$D=f\circ d_{B/A}$, donde $f:\Omega^1_{B/A}\to L$ es un
$B$-homomorfismo. En virtud de (20.5.2.6), se tiene
\[
  q_{E/B/A}(d_{B/A}(b))=d_{E/A}(q(b))
  =d_{E/A}(h(v(b)))+d_{E/A}(j(f(d_{B/A}(b))))
\]
para $b\in B$. Sea entonces
$z=\sum_i d_{B/A}(b_i)\otimes c_i$, donde $b_i\in B$ y $c_i\in C$, un
elemento de $\Omega^1_{B/A}\otimes_B C$; se tiene
\[
  (q_{E/B/A}\otimes1_C)(z)=
  \sum_i d_{E/A}(h(v(b_i)))\otimes c_i+
  \sum_i d_{E/A}(j(f(d_{B/A}(b_i))))\otimes c_i.
\]
En la primera suma, se tiene
$d_{E/A}(h(v(b_i)))=h_{E/C/A}(d_{C/A}(v(b_i)))$, luego esta suma es
$(h_{E/C/A}\otimes1_C)(v_{C/B/A}(z))$ en virtud de (20.5.2.6). Si se
toma $z\in\Upsilon_{C/B/A}$, esta suma es, por tanto, nula, y queda,
por definición de $\delta_{C/E/A}$,
\[
  (q_{E/B/A}\otimes1_C)(z)=
  \sum_i\delta_{C/E/A}(c_i f(d_{B/A}(b_i)))
  =\delta_{C/E/A}(f(z)),
\]
lo que demuestra la conmutatividad del cuadrado izquierdo de (20.6.9.1).
\end{proof}

Obsérvese que, cuando solo se supone que $\delta_{C/E/A}$ es
\emph{inyectivo} (y no necesariamente invertible por la izquierda), esta
interpretación permitiría aún definir $\chi_E$ como la
restricción de $q_{E/B/A}\otimes1_C$ a $\Upsilon_{C/B/A}$.

\begin{corollary}[20.6.10]
\label{0.20.6.10-es}
Si $\mathfrak{R}$ es un ideal de $B$, $C=B/\mathfrak{R}$,
$E=B/\mathfrak{R}^2$, y si $E$ es una $B$-extensión $A$-trivial
(18.3.7) de $C$ por $\mathfrak{R}/\mathfrak{R}^2$, entonces
el homomorfismo característico
$\chi_E:\Upsilon_{C/B/A}\to\mathfrak{R}/\mathfrak{R}^2$ es biyectivo.
\end{corollary}
\begin{proof}
\oldpage[0\textsubscript{IV-1}]{140}
En efecto, en el diagrama (20.6.9.1), las dos flechas verticales de la
derecha son homomorfismos biyectivos (20.5.12, (ii)).
\end{proof}

\begin{theorem}[20.6.11]
\label{0.20.6.11-es}
Sean $u:A\to B$, $v:B\to C$ dos homomorfismos de anillos, y $L$ un
$C$-módulo. Supongamos verificada una de las condiciones siguientes:
\begin{enumerate}
  \item[\rm(i)] $L$ es un $C$-módulo inyectivo.
  \item[\rm(ii)] $\Upsilon_{C/B/A}$ es factor directo del $C$-módulo
$\Omega^1_{B/A}\otimes_B C$ y
$u_{C/B/A}:\Omega^1_{C/A}\to\Omega^1_{C/B}$ es invertible por la derecha.
\end{enumerate}
Entonces el homomorfismo canónico (20.6.7.1)
\[
  \operatorname{Exalcom}_{B/A}(C,L)\longrightarrow
  \operatorname{Hom}_C(\Upsilon_{C/B/A},L)
\]
es biyectivo.

En particular, si $\Omega^1_{C/B}$ y $\Omega^1_{C/A}$ son
$C$-módulos proyectivos, el homomorfismo canónico (20.6.7.1) es
biyectivo.
\end{theorem}

\begin{proof}
El hecho de que cada una de las condiciones (i), (ii) implique que (20.6.7.1) es
biyectivo se sigue, en ambos casos, de la definición de $\alpha_1$. Señalemos
además que la condición (ii) es \emph{necesaria y
suficiente} para que el homomorfismo (20.6.7.1) sea biyectivo para
\emph{todo} $C$-módulo $L$ (Bourbaki, \emph{Alg.}, cap. II,
3\textsuperscript{a} ed., \S~2, n\textsuperscript{o}~1, prop.~1). Si se
supone que $\Omega^1_{C/B}$ y $\Omega^1_{C/A}$ son $C$-módulos
proyectivos, entonces, en la sucesión exacta (20.6.1.1),
$\operatorname{Ker}(u_{C/B/A})$ es un $C$-módulo proyectivo, pues la sucesión exacta
\[
  0\longrightarrow\operatorname{Ker}(u_{C/B/A})\longrightarrow\Omega^1_{C/A}
  \longrightarrow\Omega^1_{C/B}\longrightarrow0
\]
es escindida, ya que $\Omega^1_{C/B}$ es proyectivo; como
$\operatorname{Ker}(u_{C/B/A})=\operatorname{Im}(v_{C/B/A})$, la sucesión exacta
\[
  0\longrightarrow\Upsilon_{C/B/A}\longrightarrow
  \Omega^1_{B/A}\otimes_B C\longrightarrow
  \operatorname{Im}(v_{C/B/A})\longrightarrow0
\]
es escindida.
\end{proof}

\begin{corollary}[20.6.12]
\label{0.20.6.12-es}
Supongamos que $C$ es una $A$-álgebra formalmente lisa. Existe entonces
un homomorfismo canónico
\begin{equation}
\tag{20.6.12.1}
\label{0.20.6.12.1-es}
  \operatorname{Exalcom}_B(C,L)\longrightarrow
  \operatorname{Hom}_C(\Upsilon_{C/B/A},L).
\end{equation}
Además, este homomorfismo es biyectivo si se verifica una de las condiciones (i),
(ii) de (20.6.11).
\end{corollary}

\begin{proof}
En efecto, se sigue de la hipótesis sobre $C$ que
$\operatorname{Exalcom}_B(C,L)=\operatorname{Exalcom}_{B/A}(C,L)$, y
el homomorfismo (20.6.12.1) no es otro que (20.6.7.1).
\end{proof}

\begin{corollary}[20.6.13]
\label{0.20.6.13-es}
Si $C$ es una $A$-álgebra formalmente lisa y si $\Omega^1_{C/B}$ es
un $C$-módulo proyectivo, el homomorfismo (20.6.12.1) es biyectivo.
\end{corollary}

\begin{proof}
En efecto, se sabe entonces que $\Omega^1_{C/A}$ es un $C$-módulo proyectivo
(20.4.9), siendo discretas las topologías.
\end{proof}

\begin{env}[20.6.14]
\label{0.20.6.14-es}
Manteniendo las mismas notaciones, supongamos ahora además que $A$,
$B$, $C$ son $\Lambda$-álgebras y que $u$, $v$ son
$\Lambda$-homomorfismos, lo que equivale a darse tres homomorfismos
de anillos
\[
  \Lambda\xrightarrow{s}A\xrightarrow{u}B\xrightarrow{v}C.
\]

\oldpage[0\textsubscript{IV-1}]{141}
Se tienen, pues, además del módulo de imperfección $\Upsilon_{C/B/A}$, los
módulos de imperfección $\Upsilon_{B/A/\Lambda}$,
$\Upsilon_{C/A/\Lambda}$ y $\Upsilon_{C/B/\Lambda}$, y ya se han definido
homomorfismos canónicos de $C$-módulos (20.6.4.2)
\begin{align}
\tag{20.6.14.1}
\label{0.20.6.14.1-es}
  u':\Upsilon_{C/A/\Lambda}&\longrightarrow\Upsilon_{C/B/\Lambda},\\
\tag{20.6.14.2}
\label{0.20.6.14.2-es}
  s':\Upsilon_{C/B/\Lambda}&\longrightarrow\Upsilon_{C/B/A}.
\end{align}

Como en el diagrama conmutativo (20.5.7.3)
\begin{equation}
\tag{20.6.14.3}
\label{0.20.6.14.3-es}
\xymatrix{
  \Omega^1_{A/\Lambda}\otimes_A B \ar[r] \ar[d] &
  \Omega^1_{B/\Lambda} \ar[r] \ar[d] &
  \Omega^1_{B/A} \ar[r] \ar[d] & 0 \\
  \Omega^1_{A/\Lambda}\otimes_A C \ar[r] &
  \Omega^1_{C/\Lambda} \ar[r] & \Omega^1_{C/A} \ar[r] & 0
}
\end{equation}
la fila inferior está formada por $C$-módulos, se deduce por
tensorización un diagrama conmutativo
\begin{equation}
\tag{20.6.14.4}
\label{0.20.6.14.4-es}
\xymatrix{
  \Omega^1_{A/\Lambda}\otimes_A C \ar[r] \ar@{=}[d] &
  \Omega^1_{B/\Lambda}\otimes_B C \ar[r] \ar[d] &
  \Omega^1_{B/A}\otimes_B C \ar[r] \ar[d] & 0 \\
  \Omega^1_{A/\Lambda}\otimes_A C \ar[r] &
  \Omega^1_{C/\Lambda} \ar[r] & \Omega^1_{C/A} \ar[r] & 0
}
\end{equation}
donde la primera fila sigue siendo exacta y la flecha vertical de la izquierda
es la identidad; si se pone
\begin{equation}
\tag{20.6.14.5}
\label{0.20.6.14.5-es}
  \Upsilon^C_{B/A/\Lambda}=
  \operatorname{Ker}(u_{B/A/\Lambda}\otimes1_C)=
  \operatorname{Ker}(\Omega^1_{A/\Lambda}\otimes_A C\longrightarrow
       \Omega^1_{B/\Lambda}\otimes_B C)
\end{equation}
se ve, teniendo en cuenta la definición de $\Upsilon_{C/A/\Lambda}$, que
se tiene un único $C$-homomorfismo
\begin{equation}
\tag{20.6.14.6}
\label{0.20.6.14.6-es}
  v':\Upsilon^C_{B/A/\Lambda}\longrightarrow\Upsilon_{C/A/\Lambda}
\end{equation}
que hace conmutativo el diagrama
\begin{equation}
\tag{20.6.14.7}
\label{0.20.6.14.7-es}
\xymatrix{
  0 \ar[r] & \Upsilon^C_{B/A/\Lambda} \ar[r] \ar[d]^{v'} &
  \Omega^1_{A/\Lambda}\otimes_A C \ar[r] \ar@{=}[d] &
  \Omega^1_{B/\Lambda}\otimes_B C \ar[r] \ar[d] &
  \Omega^1_{B/A}\otimes_B C \ar[r] \ar[d] & 0 \\
  0 \ar[r] & \Upsilon_{C/A/\Lambda} \ar[r] &
  \Omega^1_{A/\Lambda}\otimes_A C \ar[r] &
  \Omega^1_{C/\Lambda} \ar[r] & \Omega^1_{C/A} \ar[r] & 0
}
\end{equation}
cuyas filas son exactas.

Cuando $\Lambda=\mathbb{Z}$, escribiremos $\Upsilon^C_{B/A}$ en lugar de
$\Upsilon^C_{B/A/\mathbb{Z}}$. Si $\Lambda$ es un cuerpo primo, se tiene
$\Upsilon^C_{B/A/\Lambda}=\Upsilon^C_{B/A}$.
\end{env}

\begin{env}[20.6.15]
\label{0.20.6.15-es}
Para estudiar las relaciones entre los módulos anteriores, introduciremos
por una parte el complejo de $B$-módulos
$\mathrm{K}_\bullet(B/A/\Lambda)$, y por otra, los complejos
\oldpage[0\textsubscript{IV-1}]{142}
de $C$-módulos $\mathrm{K}_\bullet(C/A/\Lambda)$ y
$\mathrm{K}_\bullet(C/B/\Lambda)$ (20.6.5), y además los complejos de
$C$-módulos siguientes. Pondremos en primer lugar
\begin{equation}
\tag{20.6.15.1}
\label{0.20.6.15.1-es}
  \mathrm{K}^C_\bullet(B/A/\Lambda)=
  \mathrm{K}_\bullet(B/A/\Lambda)\otimes_B C.
\end{equation}
Por otra parte, designaremos por $\mathrm{T}_\bullet(C/B/\Lambda)$ el
complejo de $C$-módulos cuyos términos son nulos salvo en grados 0 y 1,
y en el que
\begin{equation}
\tag{20.6.15.2}
\label{0.20.6.15.2-es}
  \mathrm{T}_0(C/B/\Lambda)=\mathrm{T}_1(C/B/\Lambda)=
  \Omega^1_{B/\Lambda}\otimes_B C
\end{equation}
siendo el operador diferencial la identidad, de modo que este complejo es
homotópico a 0; pongamos finalmente
\begin{equation}
\tag{20.6.15.3}
\label{0.20.6.15.3-es}
  \mathrm{K}'_\bullet(C/A/\Lambda)=
  \mathrm{K}_\bullet(C/A/\Lambda)\oplus
  \mathrm{T}_\bullet(C/B/\Lambda).
\end{equation}
En virtud del carácter trivial de $\mathrm{T}_\bullet$, está claro que
se tiene
\begin{equation}
\tag{20.6.15.4}
\label{0.20.6.15.4-es}
  \mathrm{H}^i(\mathrm{K}'_\bullet,L)=
  \mathrm{H}^i(\mathrm{K}_\bullet,L)
  \quad\text{y}\quad
  \mathrm{H}_i(\mathrm{K}'_\bullet,L)=
  \mathrm{H}_i(\mathrm{K}_\bullet,L)
\end{equation}
para todo $C$-módulo $L$ y todo $i$.
\end{env}
\begin{env}[20.6.16]
\label{0.20.6.16-es}
Definamos ahora una sucesión exacta de complejos, \emph{escindida} en cada
grado
\begin{equation}
\tag{20.6.16.1}
\label{0.20.6.16.1-es}
  0\to\mathrm{K}^C_\bullet(B/A/\Lambda)
  \xrightarrow{j}\mathrm{K}'_\bullet(C/A/\Lambda)
  \xrightarrow{p}\mathrm{K}_\bullet(C/B/\Lambda)\to0
\end{equation}
del modo siguiente: denotemos por un momento con
\[
  f:\Omega^1_{A/\Lambda}\otimes_A C\to
    \Omega^1_{B/\Lambda}\otimes_B C,
  \qquad
  g:\Omega^1_{B/\Lambda}\otimes_B C\to\Omega^1_{C/\Lambda}
\]
los homomorfismos canónicos $u_{B/A/\Lambda}\otimes1_C$ y
$v_{C/B/\Lambda}$ respectivamente, cuya composición es
$g\circ f=(v\circ u)_{C/A/\Lambda}$ (cf.\ \sref{0.20.6.14.4-es}).
Se tomarán $j_1(x)=(x,f(x))$, $p_1(y,z)=z-f(y)$,
$j_0(x)=(g(x),x)$, $p_0(y,z)=g(z)-y$, de modo que
$\operatorname{Im}(j_1)=\operatorname{Ker}(p_1)$ es el grafo de $f$,
complementario de $\{0\}\oplus\mathrm{T}_1$, y
$\operatorname{Im}(j_0)=\operatorname{Ker}(p_0)$ es el grafo de $g$,
complementario de $\mathrm{K}_1(C/A/\Lambda)\oplus\{0\}$; la verificación
de la conmutatividad del diagrama
\[
\xymatrix{
  0 \ar[r] & \mathrm{K}^C_1(B/A/\Lambda) \ar[r]^{j_1} \ar[d] &
  \mathrm{K}'_1(C/A/\Lambda) \ar[r]^{p_1} \ar[d] &
  \mathrm{K}_1(C/B/\Lambda) \ar[r] \ar[d] & 0 \\
  0 \ar[r] & \mathrm{K}^C_0(B/A/\Lambda) \ar[r]_{j_0} &
  \mathrm{K}'_0(C/A/\Lambda) \ar[r]_{p_0} &
  \mathrm{K}_0(C/B/\Lambda) \ar[r] & 0
}
\]
donde las flechas verticales son los operadores diferenciales, es inmediata.
\end{env}

\begin{theorem}[20.6.17]
\label{0.20.6.17-es}
Se tiene una sucesión exacta de $C$-módulos
\begin{equation}
\tag{20.6.17.1}
\label{0.20.6.17.1-es}
  0\to\Upsilon^C_{B/A/\Lambda}\xrightarrow{v'}
  \Upsilon_{C/A/\Lambda}\xrightarrow{u'}
  \Upsilon_{C/B/\Lambda}\xrightarrow{\partial}
  \Omega^1_{B/A}\otimes_B C\xrightarrow{v_{C/B/A}}
  \Omega^1_{C/A}\xrightarrow{u_{C/B/A}}\Omega^1_{C/B}\to0
\end{equation}
\oldpage[0\textsubscript{IV-1}]{143}
donde el operador de borde $\partial$ es la composición
\begin{equation}
\tag{20.6.17.2}
\label{0.20.6.17.2-es}
  \Upsilon_{C/B/\Lambda}\xrightarrow{s'}\Upsilon_{C/B/A}
  \longrightarrow\Omega^1_{B/A}\otimes_B C
\end{equation}
siendo la segunda flecha la inyección canónica.
\end{theorem}

\begin{proof}
Escribiendo la sucesión exacta de homología para la sucesión exacta de complejos
\sref{0.20.6.16.1-es}, se obtiene \sref{0.20.6.17.1-es}, pues la homología es
nula en grados distintos de 0 y de 1; el hecho de que los homomorfismos de
esta sucesión exacta que provienen por funtorialidad de $j$ y de $p$ sean
precisamente los del enunciado es inmediato. Queda por comprobar que $\partial$ es igual
a \sref{0.20.6.17.2-es}; ahora bien, un elemento $z\in\Upsilon_{C/B/\Lambda}$ es
la imagen por $p_1$ de $(0,z)$, de donde se deduce inmediatamente que $\partial(z)$ es
la imagen de $z$ por el homomorfismo canónico
\[
  s_{B/A/\Lambda}\otimes1_C:
  \Omega^1_{B/\Lambda}\otimes_B C\to\Omega^1_{B/A}\otimes_B C.
\]
La aserción se deduce de la conmutatividad del diagrama
\begin{equation}
\tag{20.6.17.3}
\label{0.20.6.17.3-es}
\xymatrix{
  \Upsilon_{C/B/\Lambda} \ar[r]^{s'} \ar[dd] & \Upsilon_{C/B/A} \ar[dd] \\
  & \\
  \Omega^1_{B/\Lambda}\otimes_B C \ar[r] &
  \Omega^1_{B/A}\otimes_B C
}
\end{equation}
(cf.\ \sref{0.20.6.4.3-es}).
\end{proof}

\begin{corollary}[20.6.18]
\label{0.20.6.18-es}
\begin{enumerate}
  \item[\rm(i)] La sucesión de $C$-módulos
\begin{equation}
\tag{20.6.18.1}
\label{0.20.6.18.1-es}
  0\to\Upsilon^C_{B/A/\Lambda}\xrightarrow{v'}
  \Upsilon_{C/A/\Lambda}\xrightarrow{u'}
  \Upsilon_{C/B/\Lambda}\xrightarrow{s'}\Upsilon_{C/B/A}\to0
\end{equation}
es exacta.
  \item[\rm(ii)] Si $B$ es un $A$-álgebra formalmente lisa relativamente
  a $\Lambda$, se tiene $\Upsilon^C_{B/A/\Lambda}=0$.
\end{enumerate}
\end{corollary}

\begin{proof}
(i) En \sref{0.20.6.17.1-es}, la imagen de $\Upsilon_{C/B/\Lambda}$ por
$\partial$ es el núcleo de $v_{C/B/A}$, es decir, $\Upsilon_{C/B/A}$ por
definición.

(ii) La hipótesis implica que la sucesión
\[
  0\to\Omega^1_{A/\Lambda}\otimes_A B\to
  \Omega^1_{B/\Lambda}\to\Omega^1_{B/A}\to0
\]
es exacta y escindida \sref{0.20.5.7-es}; por tensorización con $C$, la
sucesión
\[
  0\to\Omega^1_{A/\Lambda}\otimes_A C\to
  \Omega^1_{B/\Lambda}\otimes_B C\to
  \Omega^1_{B/A}\otimes_B C\to0
\]
permanece, pues, exacta, de donde se deduce la aserción.
\end{proof}

\begin{corollary}[20.6.19]
\label{0.20.6.19-es}
Sean $K$ un cuerpo y $E$, $F$ dos extensiones de $K$ tales que
$K\subset E\subset F$.
\begin{enumerate}
  \item[\rm(i)] Si $F$ es una extensión separable de $E$, se tiene
  $\Upsilon_{F/E/K}=0$ (es decir, el homomorfismo canónico
  $\Omega^1_{E/K}\otimes_E F\to\Omega^1_{F/K}$ es inyectivo).
  \item[\rm(ii)] Recíprocamente, si $F$ es una extensión separable de $K$ y
  si se tiene $\Upsilon_{F/E/K}=0$, entonces $F$ es una extensión separable de
  $E$.
\end{enumerate}
\end{corollary}

\begin{proof}
Si $P$ es el cuerpo primo de $K$, se tiene, en efecto, la sucesión exacta
\sref{0.20.6.18-es}
\[
  \Upsilon_{F/K/P}\to\Upsilon_{F/E/P}\to\Upsilon_{F/E/K}\to0.
\]
Si $\Upsilon_{F/E/P}=\Upsilon_{F/E}=0$, se tiene, pues,
$\Upsilon_{F/E/K}=0$, de donde (i) en virtud de \sref{0.20.6.3-es};
recíprocamente, si $\Upsilon_{F/E/K}=0$ y
$\Upsilon_{F/K/P}=\Upsilon_{F/K}=0$, se tiene
$\Upsilon_{F/E/P}=0$, de donde (ii) en virtud de \sref{0.20.6.3-es}.
\end{proof}

\oldpage[0\textsubscript{IV-1}]{144}
\begin{corollary}[20.6.20]
\label{0.20.6.20-es}
\begin{enumerate}
  \item[\rm(i)] Si $K$ es una extensión algebraica separable de $k$, se tiene
  $\Omega^1_{K/k}=0$.
  \item[\rm(ii)] Sean $k$ un cuerpo de característica 0 y $K$ una extensión
  de $k$. Para que $\Omega^1_{K/k}=0$, es necesario y suficiente que $K$ sea una
  extensión algebraica de $k$. En particular, para que un cuerpo $K$ de
  característica 0 sea tal que $\Omega^1_K=0$, es necesario y suficiente que
  $K$ sea una extensión algebraica de $\mathbb{Q}$ (cf.\ (21.4.4) y
  (21.7.5)).
\end{enumerate}
\end{corollary}

\begin{proof}
(i) Para todo $x\in K$, sea $f$ el polinomio minimal de $x$ sobre $k$; como
$f'(x)\ne0$ y $d_{K/k}(f(x))=f'(x)d_{K/k}(x)=0$, se tiene
$d_{K/k}(x)=0$ y la aserción se deduce de \sref{0.20.4.7-es}.

(ii) Existe una extensión puramente trascendente $L$ de $k$ tal que
$k\subset L\subset K$ y que $K$ es una extensión algebraica de $L$.
Como $K$ es separable sobre $L$, se deduce de \sref{0.20.6.19-es}, (i) que
la sucesión \sref{0.20.5.7.2-es}
\[
  0\to\Omega^1_{L/k}\otimes_L K\to\Omega^1_{K/k}
  \to\Omega^1_{K/L}\to0
\]
es exacta, y de (i) se deduce que $\Omega^1_{K/L}=0$. Por tanto, la relación
$\Omega^1_{K/k}=0$ equivale a $\Omega^1_{L/k}=0$, y como
$L$ es una extensión puramente trascendente de $k$, se deduce de \sref{0.20.5.10-es} que
la relación $\Omega^1_{L/k}=0$ equivale a $L=k$.
\end{proof}

\begin{remarks}[20.6.21]
\label{0.20.6.21-es}
\begin{enumerate}
  \item[\rm(i)] Como $\Upsilon_{B/A/\Lambda}$ es el núcleo de
  \[
    u_{B/A/\Lambda}:\Omega^1_{A/\Lambda}\otimes_A B
    \to\Omega^1_{B/\Lambda}
  \]
  y $\Upsilon^C_{B/A/\Lambda}$ el núcleo de
  $u_{B/A/\Lambda}\otimes1_C$, se tiene un homomorfismo canónico
\begin{equation}
\tag{20.6.21.1}
\label{0.20.6.21.1-es}
  \Upsilon_{B/A/\Lambda}\otimes_B C\to\Upsilon^C_{B/A/\Lambda}.
\end{equation}
Este homomorfismo es biyectivo cuando la sucesión
\begin{equation}
\tag{20.6.21.2}
\label{0.20.6.21.2-es}
  0\to\Upsilon_{B/A/\Lambda}\otimes_B C\to
  \Omega^1_{A/\Lambda}\otimes_A C\to
  \Omega^1_{B/\Lambda}\otimes_B C\to
  \Omega^1_{B/A}\otimes_B C\to0
\end{equation}
es exacta, lo que ocurre en los casos siguientes:
\begin{enumerate}
  \item[$1^\circ$] $C$ es un $B$-módulo \emph{plano}.
  \item[$2^\circ$] Los $B$-módulos $\Omega^1_{B/\Lambda}$ y
  $\Omega^1_{B/A}$ son \emph{planos}: en efecto, entonces ocurre lo mismo con
  $\operatorname{Ker}(\Omega^1_{B/\Lambda}\to\Omega^1_{B/A})$
  ($\mathbf{0_I}$, 6.1.2), y la sucesión \sref{0.20.6.21.2-es} es entonces
  exacta en virtud de ($\mathbf{0_I}$, 6.1.2).
\end{enumerate}
  \item[\rm(ii)] Consideremos un diagrama conmutativo de homomorfismos
  de anillos
\[
\xymatrix{
  \Lambda' \ar[r] & A' \ar[r] & B' \ar[r] & C' \\
  \Lambda \ar[u] \ar[r] & A \ar[u] \ar[r] & B \ar[u] \ar[r] & C \ar[u]
}
\]
Entonces las definiciones de \sref{0.20.6.16-es} muestran que se tiene un
diagrama conmutativo de complejos (donde las flechas verticales provienen de
\sref{0.20.6.6.4-es})
\[
\xymatrix{
  0 \ar[r] & \mathrm{K}^C_\bullet(B/A/\Lambda) \ar[r] \ar[d] &
  \mathrm{K}'_\bullet(C/A/\Lambda) \ar[r] \ar[d] &
  \mathrm{K}_\bullet(C/B/\Lambda) \ar[r] \ar[d] & 0 \\
  0 \ar[r] & \mathrm{K}^{C'}_\bullet(B'/A'/\Lambda') \ar[r] &
  \mathrm{K}'_\bullet(C'/A'/\Lambda') \ar[r] &
  \mathrm{K}_\bullet(C'/B'/\Lambda') \ar[r] & 0
}
\]
\oldpage[0\textsubscript{IV-1}]{145}
de donde se obtiene, pasando a la homología, un \emph{diagrama conmutativo}
\begin{equation}
\tag{20.6.21.3}
\label{0.20.6.21.3-es}
\xymatrix@C=0.7em{
  0 \ar[r] & \Upsilon^C_{B/A/\Lambda} \ar[r] \ar[d] &
  \Upsilon_{C/A/\Lambda} \ar[r] \ar[d] &
  \Upsilon_{C/B/\Lambda} \ar[r] \ar[d] &
  \Omega^1_{B/A}\otimes_B C \ar[r] \ar[d] &
  \Omega^1_{C/A} \ar[r] \ar[d] &
  \Omega^1_{C/B} \ar[r] \ar[d] & 0 \\
  0 \ar[r] & \Upsilon^{C'}_{B'/A'/\Lambda'} \ar[r] &
  \Upsilon_{C'/A'/\Lambda'} \ar[r] &
  \Upsilon_{C'/B'/\Lambda'} \ar[r] &
  \Omega^1_{B'/A'}\otimes_{B'} C' \ar[r] &
  \Omega^1_{C'/A'} \ar[r] & \Omega^1_{C'/B'} \ar[r] & 0
}
\end{equation}
Se tiene un diagrama conmutativo análogo para \sref{0.20.6.18.1-es}.
\end{enumerate}
\end{remarks}

\begin{proposition}[20.6.22]
\label{0.20.6.22-es}
Sean $s:\Lambda\to A$, $u:A\to B$ dos homomorfismos de anillos,
$\mathfrak{R}$ un ideal de $B$ y $C$ el anillo cociente $B/\mathfrak{R}$.
Supongamos que $E=B/\mathfrak{R}^2$ es una $B$-extensión
$\Lambda$-trivial de $C$ por $\mathfrak{R}/\mathfrak{R}^2$. Se tiene entonces una
sucesión exacta
\begin{equation}
\tag{20.6.22.1}
\label{0.20.6.22.1-es}
  0\to\Upsilon^C_{B/A/\Lambda}\xrightarrow{v'}
  \Upsilon_{C/A/\Lambda}\xrightarrow{\chi_E}
  \mathfrak{R}/\mathfrak{R}^2\xrightarrow{\delta_{C/B/A}}
  \Omega^1_{B/A}\otimes_B C\xrightarrow{v_{C/B/A}}
  \Omega^1_{C/A}\to0.
\end{equation}
\end{proposition}
 
\begin{proof}
En efecto, como $v:B\to C$ es sobreyectivo, se tiene $\Omega^1_{C/B}=0$
\sref{0.20.4.12-es}. Además, resulta de \sref{0.20.6.10-es} que
$\Upsilon_{C/B/\Lambda}$ se identifica canónicamente con
$\mathfrak{R}/\mathfrak{R}^2$. Basta entonces aplicar las sucesiones exactas
\sref{0.20.6.17.1-es} y \sref{0.20.6.18.1-es}, observando que se tiene un
diagrama conmutativo
\[
\xymatrix{
  \Upsilon_{C/A/\Lambda} \ar[r] \ar[d]_{\chi_E} &
  \Omega^1_{A/\Lambda}\otimes_A C \ar[d] \\
  \Upsilon_{C/E/\Lambda}=\Upsilon_{C/B/\Lambda}=\mathfrak{R}/\mathfrak{R}^2
    \ar[r]^-{\delta_{C/B/A}} \ar[d]_{\ell} &
  \Omega^1_{B/A}\otimes_B C=\Omega^1_{E/A}\otimes_E C \ar[d] \\
  \mathfrak{R}/\mathfrak{R}^2 \ar[r]^{\delta_{C/B/A}} &
  \Omega^1_{B/A}\otimes_B C
}
\]
y utilizando la conmutatividad del diagrama \sref{0.20.6.17.3-es}.
\end{proof}

Así se ha precisado el \emph{núcleo} de $\delta_{C/B/A}$ en el caso en que
existe un anillo $\Lambda$ tal que $A$ sea una $\Lambda$-álgebra y que
$B/\mathfrak{R}^2$ sea $\Lambda$-trivial; obsérvese que este es el caso
en particular cuando $C$ es una \emph{$\Lambda$-álgebra formalmente
lisa} (para la topología discreta).

Supongamos además que se tiene un diagrama conmutativo de homomorfismos
de anillos
\[
\xymatrix{
  \Lambda' \ar[r] & A' \ar[r] & B' \\
  \Lambda \ar[u] \ar[r] & A \ar[u] \ar[r] & B \ar[u]^f
}
\]
\oldpage[0\textsubscript{IV-1}]{146}
que $\mathfrak{R}'$ es un ideal de $B'$ tal que
$f(\mathfrak{R})\subset\mathfrak{R}'$, y que
$E'=B'/\mathfrak{R}'^2$ sea una $B'$-extensión $\Lambda'$-trivial de
$C'=B'/\mathfrak{R}'$ por $\mathfrak{R}'/\mathfrak{R}'^2$. Se tiene entonces un
diagrama conmutativo
\begin{equation}
\tag{20.6.22.2}
\label{0.20.6.22.2-es}
\xymatrix@C=0.7em{
  0 \ar[r] & \Upsilon^{C'}_{B'/A'/\Lambda'} \ar[r] &
  \Upsilon_{C'/A'/\Lambda'} \ar[r] &
  \mathfrak{R}'/\mathfrak{R}'^2 \ar[r] &
  \Omega^1_{B'/A'}\otimes_{B'} C' \ar[r] &
  \Omega^1_{C'/A'} \ar[r] & 0 \\
  0 \ar[r] & \Upsilon^C_{B/A/\Lambda} \ar[r] \ar[u] &
  \Upsilon_{C/A/\Lambda} \ar[r] \ar[u] &
  \mathfrak{R}/\mathfrak{R}^2 \ar[r] \ar[u] &
  \Omega^1_{B/A}\otimes_B C \ar[r] \ar[u] &
  \Omega^1_{C/A} \ar[r] \ar[u] & 0
}
\end{equation}
como resulta de \sref{0.20.6.21.3-es} y \sref{0.20.5.11.3-es}.

\begin{corollary}[20.6.23]
\label{0.20.6.23-es}
Bajo las hipótesis de \sref{0.20.6.22-es}, supongamos además que $B$ sea
una $A$-álgebra formalmente lisa. Se tiene entonces una sucesión exacta
\begin{equation}
\tag{20.6.23.1}
\label{0.20.6.23.1-es}
  0\to\Upsilon_{C/A/\Lambda}\xrightarrow{\chi_E}
  \mathfrak{R}/\mathfrak{R}^2\xrightarrow{\delta_{C/B/A}}
  \Omega^1_{B/A}\otimes_B C\to\Omega^1_{C/A}\to0.
\end{equation}
\end{corollary}

\begin{proof}
Esto resulta, en efecto, de \sref{0.20.6.18-es}, (ii).
\end{proof}

\begin{env}[20.6.24]
\label{0.20.6.24-es}
Cuando se satisfacen las hipótesis de \sref{0.20.6.22-es}, se dice también
que el homomorfismo característico $\chi_E$ es el
\emph{homomorfismo característico de la $A$-álgebra $B$, respecto de
$\Lambda$ y del ideal $\mathfrak{R}$} (omitiéndose estas últimas
precisiones cuando no haya riesgo de confusión); se denotará
a veces por $\chi_B$ o $\chi_{B/A}$.
\end{env}

\begin{proposition}[20.6.25]
\label{0.20.6.25-es}
Sean $s:\Lambda\to A$, $u:A\to B$, $v:B\to C$ tres homomorfismos
de anillos, $L$ un $C$-módulo. Se tiene entonces una sucesión exacta
\begin{equation}
\tag{20.6.25.1}
\label{0.20.6.25.1-es}
\begin{split}
  0\to\operatorname{D\acute{e}r}_B(C,L)\xrightarrow{u^1}
  \operatorname{D\acute{e}r}_A(C,L)\xrightarrow{v^1}
  \operatorname{D\acute{e}r}_A(B,L)\xrightarrow{\partial} \\
  \to\operatorname{Exalcom}_{B/\Lambda}(C,L)\xrightarrow{u^1}
  \operatorname{Exalcom}_{A/\Lambda}(C,L)\xrightarrow{v^1}
  \operatorname{Exalcom}_{A/\Lambda}(B,L)\to0
\end{split}
\end{equation}
donde $u^1$, $v^1$ son los homomorfismos definidos en (18.3.6.4) y
(18.3.6.2), y $\partial$ se define como en \sref{0.20.2.2-es}.
\end{proposition}

\begin{proof}
En efecto, como la sucesión exacta \sref{0.20.6.16.1-es} es escindida, se
deduce de ella una sucesión exacta
\[
  0\to\operatorname{Hom}_C(\mathrm{K}_\bullet(C/B/\Lambda),L)\to
  \operatorname{Hom}_C(\mathrm{K}'_\bullet(C/A/\Lambda),L)\to
  \operatorname{Hom}_C(\mathrm{K}^C_\bullet(B/A/\Lambda),L)\to0.
\]
Si se aplica a este complejo la sucesión exacta de cohomología, teniendo en cuenta
\sref{0.20.6.15.4-es} y \sref{0.20.6.6-es}, se obtiene
\sref{0.20.6.25.1-es}, pues se tiene
\[
  \operatorname{Hom}_C(\mathrm{K}^C_\bullet(B/A/\Lambda),L)
  =\operatorname{Hom}_B(\mathrm{K}_\bullet(B/A/\Lambda),L)
\]
por definición; la identificación de $u^1$ y $v^1$ con los homomorfismos
de (18.3.4.2) resulta de \sref{0.20.6.6.4-es}.
\end{proof}

\begin{remark}[20.6.26]
\label{0.20.6.26-es}
En este número, los complejos $\mathrm{K}_\bullet(C/B/A)$ han aparecido como
un artificio técnico destinado a simplificar la exposición de ciertos
comportamientos funtoriales. En realidad, estos complejos, considerados como objetos
de la categoría de los
\oldpage[0\textsubscript{IV-1}]{147}
complejos de $C$-módulos «salvo homotopía» (es decir, donde los morfismos
son las clases de homomorfismos homotópicos), son invariantes notables,
más finos que el par formado por $\Omega^1_{C/B}$ y
$\Upsilon_{C/B/A}$. Cuando $k$ es un cuerpo primo, y $C$ una $k$-álgebra
formalmente lisa (por ejemplo, un anillo regular de tipo finito sobre una
extensión de $k$ (cf.\ ($\mathbf{IV}$, 6.8.6))), se puede demostrar que el
complejo $\mathrm{K}_\bullet(C/A/k)$ puede describirse haciendo intervenir
únicamente $C$ y $A$ (con exclusión de $k$): se expresa $C$ como cociente
de un álgebra de polinomios $B$ sobre $A$ por un ideal $\mathfrak{Q}$, y se
considera el complejo $\mathrm{F}_\bullet(C/A)$ de 2 términos no nulos
\[
  \cdots\to0\to\mathfrak{Q}/\mathfrak{Q}^2\to
  \Omega^1_{B/A}\otimes_B C\to0\to\cdots
\]
(cuya homología coincide efectivamente con la de
$\mathrm{K}_\bullet(C/A/k)$ en virtud de \sref{0.20.6.23.1-es}). Estos
complejos $\mathrm{F}_\bullet(C/A)$, que desde el punto de vista del álgebra
homológica desempeñan el papel de un fibrado conormal para $\operatorname{Spec}(C)$
sobre $\operatorname{Spec}(A)$, ocuparán un lugar importante en
los capítulos de esta obra dedicados a la dualidad de haces coherentes y
al teorema de Riemann--Roch.
\end{remark}

\subsection{Generalizaciones a los anillos topológicos}
\label{subsection:0.20.7-es}

\begin{env}[20.7.1]
\label{0.20.7.1-es}
Resulta inmediatamente de las definiciones que, si en \sref{0.20.5.2-es} y
\sref{0.20.5.3-es} se supone que los anillos $A$, $B$, $C$ son anillos
topológicos y que los homomorfismos de anillos $u$, $v$ son continuos, entonces los
homomorfismos $u_{C/B/A}$ y $v_{C/B/A}$ son \emph{continuos} (naturalmente,
hay que tomar sobre $\Omega^1_{B/A}\otimes_B C$ la topología producto
tensorial).

Además, $u_{C/B/A}:\Omega^1_{C/A}\to\Omega^1_{C/B}$ es un
\emph{morfismo estricto sobreyectivo} de $C$-módulos topológicos; en efecto, lo
mismo ocurre con el homomorfismo canónico $C\otimes_A C\to C\otimes_B C$, teniendo
en cuenta la definición de la topología producto tensorial; pasando a los
cocientes se deduce que el homomorfismo correspondiente
$\mathrm{P}^1_{C/A}\to\mathrm{P}^1_{C/B}$ es un morfismo estricto sobreyectivo,
y $u_{C/B/A}$ es la restricción de este último a $\Omega^1_{C/A}$.

En \sref{0.20.5.5-es}, si $A'$ y $B$ son $A$-álgebras topológicas,
y si $B'$ y $\Omega^1_{B/A}\otimes_B B'$ están provistos de las topologías producto
tensorial, el homomorfismo canónico \sref{0.20.5.5.1-es} es un isomorfismo
topológico, teniendo en cuenta que $\mathrm{P}^1_{B/A}$ es suma directa
topológica de $B$ y de $\Omega^1_{B/A}$ \sref{0.20.4.8-es}.
\end{env}

\begin{proposition}[20.7.2]
\label{0.20.7.2-es}
Sean $u:A\to B$, $v:B\to C$ dos homomorfismos continuos de anillos
topológicos. Para que el homomorfismo continuo
$v_{C/B/A}:\Omega^1_{B/A}\otimes_B C\to\Omega^1_{C/A}$ sea formalmente
invertible por la izquierda (cf.\ \sref{0.19.1.5-es}), es necesario y suficiente que $C$
sea una $B$-álgebra formalmente lisa respecto de $A$
\sref{0.19.9.1-es} (y \emph{a fortiori} basta que $C$ sea una
$B$-álgebra formalmente lisa).
\end{proposition}

\begin{proof}
Decir que $v_{C/B/A}$ es formalmente invertible por la izquierda significa, en efecto,
puesto que las topologías de $\Omega^1_{C/A}$ y $\Omega^1_{B/A}\otimes_B C$
son menos finas que las deducidas de la topología de $C$
\sref{0.20.4.5-es}, que para todo $C$-módulo \emph{discreto} $L$, anulado por un
ideal abierto de $C$, el homomorfismo canónico
\[
  \operatorname{Hom.\ cont}_C(\Omega^1_{C/A},L)\to
  \operatorname{Hom.\ cont}_C(\Omega^1_{B/A}\otimes_B C,L)
\]
\oldpage[0\textsubscript{IV-1}]{148}
es \emph{sobreyectivo} \sref{0.19.1.5-es}; como
$\operatorname{Hom.\ cont}_C(\Omega^1_{B/A}\otimes_B C,L)=
\operatorname{Hom.\ cont}_B(\Omega^1_{B/A},L)$ por definición de la topología
producto tensorial, equivale, en virtud de \sref{0.20.4.8-es}, a decir
que el homomorfismo canónico
\[
  \operatorname{D\acute{e}r.\ cont}_A(C,L)\to
  \operatorname{D\acute{e}r.\ cont}_A(B,L)
\]
es sobreyectivo. Pero la sucesión exacta \sref{0.20.3.8.2-es} muestra que esta
condición equivale a $\operatorname{Exalcotop}_{B/A}(C,L)=0$, es decir,
precisamente a que $C$ sea formalmente lisa respecto de $A$
\sref{0.19.9.8-es}.
\end{proof}

\begin{corollary}[20.7.3]
\label{0.20.7.3-es}
Supongamos que, en $B$ y en $C$, el cuadrado de todo ideal abierto es abierto.
Para que $C$ sea una $B$-álgebra formalmente lisa respecto de $A$, es
necesario y suficiente que, si se denota por $(\mathfrak{R}_\lambda)$ un sistema
fundamental de entornos de $0$ formado por ideales de $C$, entonces, para todo
$\lambda$, el homomorfismo
\begin{equation}
\tag{20.7.3.1}
\label{0.20.7.3.1-es}
  v_{C/B/A}\otimes 1_{C/\mathfrak{R}_\lambda}:
  \Omega^1_{B/A}\otimes_B(C/\mathfrak{R}_\lambda)\to
  \Omega^1_{C/A}\otimes_C(C/\mathfrak{R}_\lambda)
\end{equation}
sea invertible por la izquierda.
\end{corollary}

\begin{proof}
Se sabe, en efecto, que la topología de $\Omega^1_{B/A}$ (resp.\
$\Omega^1_{C/A}$) se deduce de la de $B$ (resp.\ de $C$)
\sref{0.20.4.5-es}; se concluye inmediatamente que la topología de
$\Omega^1_{B/A}\otimes_B C$ se deduce también de la de $C$; el corolario
resulta entonces de \sref{0.20.7.2-es} y \sref{0.19.1.7-es}.
\end{proof}

\begin{proposition}[20.7.4]
\label{0.20.7.4-es}
Sean $A$ un anillo topológico, $B$ una $A$-álgebra topológica. Para que
$B$ sea formalmente no ramificada \sref{0.19.10.2-es}, es necesario y suficiente
que el completado separado $\widehat{\Omega}^1_{B/A}$ sea nulo.
\end{proposition}

\begin{proof}
En efecto, se deduce inmediatamente de \sref{0.19.10.4-es} y
\sref{0.20.1.1-es} que, para que $B$ sea formalmente no ramificada, es necesario y
suficiente que, para todo ideal abierto $\mathfrak{R}$ de $B$, todo ideal abierto
$\mathfrak{H}$ de $A$ tal que $\mathfrak{H}B\subset\mathfrak{R}$ y todo
$(B/\mathfrak{R})$-módulo $L$, se tenga
$\operatorname{D\acute{e}r}_{A/\mathfrak{H}}(B/\mathfrak{R},L)=0$,
es decir, $\operatorname{D\acute{e}r.\ cont}_A(B,L)=0$ para todo
$B$-módulo discreto $L$ anulado por un ideal abierto de $B$; en virtud de
\sref{0.20.4.8-es}, esto equivale a
$\operatorname{Hom.\ cont}_B(\Omega^1_{B/A},L)=0$ para tal $B$-módulo,
de donde se sigue inmediatamente la proposición.

Cuando $B$ es \emph{discreto}, la condición del enunciado de
\sref{0.20.7.4-es} equivale por tanto a $\Omega^1_{B/A}=0$.
\end{proof}

\begin{corollary}[20.7.5]
\label{0.20.7.5-es}
Sean $A$ un anillo, $\mathfrak{m}$ un ideal de $A$, $B$ una $A$-álgebra;
dotemos $A$ de la topología $\mathfrak{m}$-preádica, $B$ de la topología
deducida de la de $A$, y pongamos $A_0=A/\mathfrak{m}$,
$B_0=B/\mathfrak{m}B=B\otimes_A A_0$. Entonces, para que $B$ sea una
$A$-álgebra formalmente no ramificada, es necesario y suficiente que
$\Omega^1_{B_0/A_0}=0$ (o también que $B_0$ sea una $A_0$-álgebra
formalmente no ramificada).
\end{corollary}

\begin{proof}
En efecto, se tiene
$\Omega^1_{B/A}\otimes_B B_0=\Omega^1_{B/A}/\mathfrak{m}.\Omega^1_{B/A}
=\Omega^1_{B_0/A_0}$ en virtud de \sref{0.20.5.5-es}; escribir que todo
submódulo abierto de $\Omega^1_{B/A}$ es igual a $\Omega^1_{B/A}$ equivale
por \sref{0.20.4.5-es} a escribir que
$\mathfrak{m}.\Omega^1_{B/A}=\Omega^1_{B/A}$, de donde se obtiene la conclusión.
\end{proof}

\begin{proposition}[20.7.6]
\label{0.20.7.6-es}
Sean $u:A\to B$, $v:B\to C$ dos homomorfismos continuos de anillos
topológicos, y supongamos que $v$ haga de $C$ una $B$-álgebra formalmente
étale; entonces $v_{C/B/A}:\Omega^1_{B/A}\otimes_B C\to\Omega^1_{C/A}$ es un
bimorfismo formal \sref{0.19.1.2-es}.
\end{proposition}

\begin{proof}
En efecto, para todo $C$-módulo discreto $L$ anulado por un ideal abierto de $C$,
se tiene entonces por \sref{0.20.7.4-es}
$\operatorname{D\acute{e}r.\ cont}_B(C,L)=0$ y, por consiguiente,
\sref{0.20.3.6.1-es}, el homomorfismo canónico
\oldpage[0\textsubscript{IV-1}]{149}
$\operatorname{D\acute{e}r.\ cont}_A(C,L)\to
\operatorname{D\acute{e}r.\ cont}_A(B,L)$ es inyectivo; además, es
sobreyectivo en virtud de \sref{0.20.7.2-es}, por lo que es biyectivo; se deduce
que la imagen de $v_{C/B/A}$ es necesariamente densa en $\Omega^1_{C/A}$,
pues, de lo contrario, el cociente de $\Omega^1_{C/A}$ por la adherencia de
$\operatorname{Im}(v_{C/B/A})$ sería separado y $\neq0$ y tendría por tanto un
cociente discreto $L\neq0$, en contradicción con lo que acabamos de ver
(teniendo en cuenta \sref{0.20.4.8-es}). Por consiguiente, $v_{C/B/A}$, que es un
monomorfismo formal en virtud de \sref{0.20.7.2-es}, es también un epimorfismo formal
\sref{0.19.1.2-es}, y por tanto un bimorfismo formal.
\end{proof}

\begin{corollary}[20.7.7]
\label{0.20.7.7-es}
Supongamos que en $B$ y en $C$ el cuadrado de todo ideal abierto sea abierto.
Si $C$ es una $B$-álgebra formalmente étale, entonces, para todo ideal abierto
$\mathfrak{R}_\lambda$ de $C$, el homomorfismo \sref{0.20.7.3.1-es} es
biyectivo.
\end{corollary}

\begin{proof}
En efecto, en virtud de \sref{0.19.1.1-es}, este homomorfismo es sobreyectivo, y
es inyectivo por \sref{0.20.7.3-es}.
\end{proof}

\begin{proposition}[20.7.8]
\label{0.20.7.8-es}
Sean $u:A\to B$ un homomorfismo continuo de anillos topológicos,
$\mathfrak{R}$ un ideal de $B$, $C$ el anillo topológico cociente
$B/\mathfrak{R}$, $v:B\to C$ el homomorfismo canónico. Entonces:
\begin{enumerate}
  \item[\rm(i)] En la sucesión exacta \sref{0.20.5.12.1-es}
\[
  \mathfrak{R}/\mathfrak{R}^2\xrightarrow{\delta_{C/B/A}}
  \Omega^1_{B/A}\otimes_B C\xrightarrow{v_{C/B/A}}\Omega^1_{C/A}\to0
\]
el homomorfismo $\delta_{C/B/A}$ es continuo y el homomorfismo $v_{C/B/A}$
es un morfismo estricto de $C$-módulos topológicos.
  \item[\rm(ii)] Para que $\delta_{C/B/A}$ sea formalmente invertible por la
izquierda \sref{0.19.1.5-es}, es necesario y suficiente que, para todo $C$-módulo
discreto $L$ anulado por un ideal abierto de $C$, el homomorfismo canónico
\begin{equation}
\tag{20.7.8.1}
\label{0.20.7.8.1-es}
  \operatorname{Exalcotop}_A(C,L)\to\operatorname{Exalcotop}_A(B,L)
\end{equation}
sea inyectivo.
\end{enumerate}
\end{proposition}

\begin{proof}
(i) La primera afirmación es evidente. Por otra parte, el homomorfismo
canónico $v\otimes v:B\otimes_A B\to C\otimes_A C$ es un morfismo estricto
por definición de la topología producto tensorial, y se deduce inmediatamente
(cf.\ \sref{0.20.5.2-es}) que lo mismo ocurre con $v_{C/B/A}$.

(ii) Decir que $\delta_{C/B/A}$ es formalmente invertible por la izquierda significa
que, para todo $C$-módulo discreto $L$ anulado por un ideal abierto de $C$,
el homomorfismo canónico
\[
  \operatorname{Hom.\ cont}_C(\Omega^1_{B/A}\otimes_B C,L)\to
  \operatorname{Hom.\ cont}_C(\mathfrak{R}/\mathfrak{R}^2,L)
\]
es sobreyectivo. Ahora bien, teniendo en cuenta \sref{0.18.4.3-es} y
\sref{0.20.4.8-es}, esto equivale a decir que el homomorfismo canónico
\[
  \operatorname{D\acute{e}r.\ cont}_A(B,L)\to
  \operatorname{Exalcotop}_B(C,L)
\]
es sobreyectivo, y la conclusión se deduce entonces de la sucesión exacta
\sref{0.20.3.6.1-es}.

El hecho de que \sref{0.20.7.8.1-es} sea inyectivo se expresa aún de la manera
siguiente, teniendo en cuenta la definición de ambos miembros
\sref{0.18.4.1-es}: para todo ideal abierto $\mathfrak{M}$ de $A$, todo ideal
abierto $\mathfrak{N}$ de $B$ tal que $\mathfrak{M}B\subset\mathfrak{N}$, y
toda $(A/\mathfrak{M})$-extensión $E$ de
$B/(\mathfrak{R}+\mathfrak{N})$ por un
$(B/(\mathfrak{R}+\mathfrak{N}))$-módulo $L$, tal que la imagen recíproca de
$E$ por el homomorfismo canónico
$B/\mathfrak{N}\to B/(\mathfrak{R}+\mathfrak{N})$ sea
$(A/\mathfrak{M})$-trivial, existe un ideal abierto
$\mathfrak{M}'\subset\mathfrak{M}$ de $A$,
\oldpage[0\textsubscript{IV-1}]{150}
un ideal abierto $\mathfrak{N}'\subset\mathfrak{N}$ de $B$ tales que
$\mathfrak{M}'B\subset\mathfrak{N}'$ y que la imagen recíproca de $E$ por
el homomorfismo canónico
$B/(\mathfrak{R}+\mathfrak{N}')\to B/(\mathfrak{R}+\mathfrak{N})$ sea
$(A/\mathfrak{M}')$-trivial.
\end{proof}

\begin{corollary}[20.7.9]
\label{0.20.7.9-es}
Si la $A$-álgebra topológica $C=B/\mathfrak{R}$ es formalmente lisa,
el homomorfismo canónico $\delta_{C/B/A}$ es formalmente invertible por la
izquierda.
\end{corollary}

\begin{env}[20.7.10]
\label{0.20.7.10-es}
En \sref{0.20.6.1-es}, cuando $A$, $B$, $C$ son anillos topológicos,
$u$ y $v$ homomorfismos continuos, se dota a $\Upsilon_{C/B/A}$ de la
topología inducida por la de $\Omega^1_{B/A}\otimes_B C$; los
homomorfismos \sref{0.20.6.4.2-es} y \sref{0.20.6.4.4-es} son entonces
continuos siempre que también lo sean los del diagrama
\sref{0.20.6.4.1-es}. Además, si, en \sref{0.20.6.7-es}, se supone que
$L$ es un $C$-módulo discreto anulado por un ideal abierto de $C$, se deduce,
por paso al límite inductivo, un $C$-homomorfismo canónico
\begin{equation}
\tag{20.7.10.1}
\label{0.20.7.10.1-es}
  \operatorname{Exalcotop}_{B/A}(C,L)\to
  \operatorname{Hom.\ cont}_C(\Upsilon_{C/B/A},L)
\end{equation}
teniendo en cuenta \sref{0.18.5.3.1-es}: para todo ideal abierto
$\mathfrak{M}$ de $A$, todo ideal abierto $\mathfrak{N}$ de $B$ tal que
$\mathfrak{M}B\subset\mathfrak{N}$, todo ideal abierto $\mathfrak{P}$ de $C$
tal que $\mathfrak{N}C\subset\mathfrak{P}$ y $\mathfrak{P}.L=0$, toda
$(B/\mathfrak{N})$-extensión $E$ de $C/\mathfrak{P}$ por $L$ que sea
$(A/\mathfrak{M})$-trivial procede de los datos de un $C$-homomorfismo
continuo $f$ de $\Omega^1_{B/A}\otimes_B C$ en $L$, y el homomorfismo
\sref{0.20.7.10.1-es} hace corresponder a la imagen en
$\operatorname{Exalcotop}_{B/A}(C,L)$ de la clase de $E$, la restricción de
$f$ a $\Upsilon_{C/B/A}$, llamado \emph{homomorfismo característico de $E$}
y que también se denota $\chi_E$.
\end{env}

\begin{proposition}[20.7.11]
\label{0.20.7.11-es}
Supongamos que la topología de $C$ sea tal que el cuadrado de un ideal abierto
sea abierto. Si $C$ es una $A$-álgebra topológica formalmente lisa y si
$\Omega^1_{C/B}$ es un $C$-módulo formalmente proyectivo, se tiene un
isomorfismo canónico
\begin{equation}
\tag{20.7.11.1}
\label{0.20.7.11.1-es}
  \operatorname{Exalcotop}_B(C,L)\xrightarrow{\sim}
  \operatorname{Hom.\ cont}_C(\Upsilon_{C/B/A},L)
\end{equation}
para todo $C$-módulo discreto $L$ anulado por un ideal abierto de $C$.
\end{proposition}

\begin{proof}
En efecto, en la sucesión exacta \sref{0.20.3.7.1-es},
$\operatorname{Exalcotop}_A(C,L)=0$ \sref{0.19.4.4-es}, por lo que
$\operatorname{Exalcotop}_B(C,L)=\operatorname{Exalcotop}_{B/A}(C,L)$ y
el homomorfismo \sref{0.20.7.8.1-es} no es otro que
\sref{0.20.7.10.1-es}; el hecho de que sea biyectivo se deduce de
\sref{0.20.6.13-es} por paso al límite inductivo, teniendo en cuenta que
la topología de $\Omega^1_{C/B}$ se deduce entonces de la de $C$
\sref{0.20.4.5-es} y de \sref{0.19.2.4-es}.
\end{proof}

\begin{env}[20.7.12]
\label{0.20.7.12-es}
En \sref{0.20.6.14-es} se dota de nuevo a $\Upsilon^C_{B/A/\Lambda}$ de la
topología inducida por la de $\Omega^1_{A/\Lambda}\otimes_A C$ y entonces
el homomorfismo \sref{0.20.6.14.6-es} es continuo, cuando los anillos
considerados son topológicos y los homomorfismos de anillos son continuos.
\end{env}

\begin{env}[20.7.13]
\label{0.20.7.13-es}
Si, en \sref{0.20.6.23-es}, se supone que los anillos $\Lambda$, $A$, $B$,
$C$ son topológicos, los homomorfismos $s$, $u$, $v$ continuos y $L$ un
$C$-módulo discreto anulado por un ideal abierto de $C$, entonces se puede pasar
al límite inductivo como en \sref{0.20.3.6-es}, y se tiene una sucesión exacta
\begin{equation}
\tag{20.7.13.1}
\label{0.20.7.13.1-es}
\begin{split}
  0\to\operatorname{D\acute{e}r.\ cont}_B(C,L)\to
  \operatorname{D\acute{e}r.\ cont}_A(C,L)\to
  \operatorname{D\acute{e}r.\ cont}_A(B,L)\to \\
  \to\operatorname{Exalcotop}_{B/\Lambda}(C,L)\to
  \operatorname{Exalcotop}_{A/\Lambda}(C,L)\to
  \operatorname{Exalcotop}_{A/\Lambda}(B,L)\to0.
\end{split}
\end{equation}
\end{env}

\begin{env}[20.7.14]
\label{0.20.7.14-es}
Sean $A$ un anillo topológico, $B$ una $A$-álgebra topológica,
$\mathfrak{M}'$ un ideal abierto de $A$, $\mathfrak{N}'$ un ideal abierto de
$B$ tales que $\mathfrak{M}'B\subset\mathfrak{N}'$; pongamos
$A'=A/\mathfrak{M}'$,
\oldpage[0\textsubscript{IV-1}]{151}
$B'=B/\mathfrak{N}'$; el núcleo del homomorfismo
$B\otimes_A B\to B'\otimes_{A'}B'$ es
$\mathfrak{U}'=\operatorname{Im}(\mathfrak{N}'\otimes B+B\otimes\mathfrak{N}')$,
de donde se deduce inmediatamente que el núcleo del homomorfismo
\begin{equation}
\tag{20.7.14.1}
\label{0.20.7.14.1-es}
  \varphi_{(\mathfrak{M}',\mathfrak{N}')}:
  \Omega^1_{B/A}\to\Omega^1_{B'/A'}
\end{equation}
es $((\mathfrak{J}\cap\mathfrak{U}')+\mathfrak{J}^2)/\mathfrak{J}^2$ ; por
otra parte, el homomorfismo \sref{0.20.7.14.1-es} es sobreyectivo, como se
desprende de \sref{0.20.4.7-es}. Si $\mathfrak{M}''$ (resp.\
$\mathfrak{N}''$) es un segundo ideal abierto de $A$ (resp.\ $B$) tal que
$\mathfrak{M}''\subset\mathfrak{M}'$, $\mathfrak{N}''\subset\mathfrak{N}'$
y $\mathfrak{M}''B\subset\mathfrak{N}''$, y si se pone
$A''=A/\mathfrak{M}''$, $B''=B/\mathfrak{N}''$, se tiene de igual modo un
homomorfismo sobreyectivo
\[
  \varphi_{(\mathfrak{M}'',\mathfrak{N}''),(\mathfrak{M}',\mathfrak{N}')}:
  \Omega^1_{B''/A''}\to\Omega^1_{B'/A'}
\]
y estos homomorfismos forman evidentemente un \emph{sistema proyectivo}. Si
se observa que los
$((\mathfrak{J}\cap\mathfrak{U}')+\mathfrak{J}^2)/\mathfrak{J}^2$ forman
un sistema fundamental de entornos de $0$ en $\Omega^1_{B/A}$, se ve
entonces que el \emph{completado separado} $\widehat{\Omega}^1_{B/A}$ del
módulo $B$-topológico $\Omega^1_{B/A}$ viene dado, salvo un isomorfismo
canónico, por
\begin{equation}
\tag{20.7.14.2}
\label{0.20.7.14.2-es}
  \widehat{\Omega}^1_{B/A}=\varprojlim(\Omega^1_{B'/A'}).
\end{equation}

Además, el homomorfismo canónico
$j:\Omega^1_{B/A}\to\widehat{\Omega}^1_{B/A}$ es el límite proyectivo del
sistema proyectivo de los $\varphi_{(\mathfrak{M}',\mathfrak{N}')}$, por lo que
$\varphi_{(\mathfrak{M}',\mathfrak{N}')}$ se factoriza como
$\Omega^1_{B/A}\xrightarrow{j}\widehat{\Omega}^1_{B/A}\to
\Omega^1_{B'/A'}$, y como es sobreyectivo se concluye que el homomorfismo
canónico
\begin{equation}
\tag{20.7.14.3}
\label{0.20.7.14.3-es}
  \widehat{\Omega}^1_{B/A}\to\Omega^1_{B'/A'}
\end{equation}
es sobreyectivo para todo par $(\mathfrak{M}',\mathfrak{N}')$.

Además, en lo anterior se puede reemplazar en todas partes $A'$ por $A$.

Finalmente, si $L$ es un $B$-módulo topológico \emph{separado y completo}, todo
$B$-homomorfismo continuo de $\Omega^1_{B/A}$ en $L$ se prolonga de forma
única a un $\widehat B$-homomorfismo continuo de
$\widehat{\Omega}^1_{B/A}$ en $L$, y recíprocamente un homomorfismo de este
tipo vuelve a dar por composición con $\Omega^1_{B/A}\to\widehat{\Omega}^1_{B/A}$
un $B$-homomorfismo continuo, por lo que se tiene un isomorfismo canónico
\[
  \operatorname{Hom.\ cont}_B(\Omega^1_{B/A},L)\xrightarrow{\sim}
  \operatorname{Hom.\ cont}_{\widehat B}(\widehat{\Omega}^1_{B/A},L).
\]
En particular, si $L$ es un $B$-módulo discreto anulado por un ideal
abierto de $B$, se ve que el isomorfismo canónico
\sref{0.20.4.8.2-es} también puede escribirse
\begin{equation}
\tag{20.7.14.4}
\label{0.20.7.14.4-es}
  \operatorname{Hom.\ cont}_{\widehat B}(\widehat{\Omega}^1_{B/A},L)
  \xrightarrow{\sim}\operatorname{D\acute{e}r.\ cont}_A(B,L).
\end{equation}
\end{env}

\begin{proposition}[20.7.15]
\label{0.20.7.15-es}
Sean $A$ un anillo discreto, $B$ un $A$-álgebra topológica adica
($\mathbf{0}_{\mathrm I}$, 7.1.9), $\mathfrak{n}$ un ideal de definición de
$B$, $B_0=B/\mathfrak{n}$. Se supone que $\Omega^1_{B_0/A}$ y
$\mathfrak{n}/\mathfrak{n}^2$ son módulos de tipo finito sobre $B_0$. Entonces
$\widehat{\Omega}^1_{B/A}$ es un $\widehat B$-módulo de tipo finito.
\end{proposition}

\begin{proof}
Como el cuadrado de todo ideal abierto de $B$ es abierto, la topología de
$\Omega^1_{B/A}$ es la topología $\mathfrak{n}$-preádica
\sref{0.20.4.5-es}. Teniendo en cuenta la hipótesis de que $B$ es un anillo adico,
basta, por tanto, en virtud de ($\mathbf{0}_{\mathrm I}$, 7.2.7 y 7.2.9), demostrar
\oldpage[0\textsubscript{IV-1}]{152}
que $\Omega^1_{B/A}/\mathfrak{n}.\Omega^1_{B/A}=
\Omega^1_{B/A}\otimes_B B_0$ es un $B_0$-módulo de tipo finito. Pero esto
se deduce de la hipótesis y de la sucesión exacta \sref{0.20.5.12.1-es}
\[
  \mathfrak{n}/\mathfrak{n}^2\to
  \Omega^1_{B/A}\otimes_B B_0\to\Omega^1_{B_0/A}\to0.
\]
\end{proof}

\begin{env}[20.7.16]
\label{0.20.7.16-es}
La proposición \sref{0.20.7.15-es} se aplica, por ejemplo, cuando $A$ es un
cuerpo $k$, $B=k'[[T_1,\ldots,T_n]]$ es un álgebra de series formales provista de
su topología usual, y $k'$ es una extensión \emph{finita} de $k$ (cf.\
\sref{0.21.9.2-es}). Se observará que el cuerpo de fracciones $K$ de $B$ tiene
grado de trascendencia infinito sobre $k$; cuando $k$ tiene característica
$0$, se deduce inmediatamente (con ayuda, en particular, de
\sref{0.20.4.13-es}, (i), y de \sref{0.20.5.9-es}, usando también el hecho de
que toda derivación de un cuerpo de característica $0$ se prolonga a toda
extensión) que $\Omega^1_{B/k}$ \emph{no es un $B$-módulo de tipo finito}.
\end{env}

\begin{env}[20.7.17]
\label{0.20.7.17-es}
Sean $A$, $B$, $C$ tres anillos topológicos, y sean $u:A\to B$, $v:B\to C$
dos homomorfismos continuos; reemplazando $A$, $B$, $C$ por cocientes por
ideales abiertos $A'=A/\mathfrak{M}'$, $B'=B/\mathfrak{N}'$,
$C'=C/\mathfrak{P}'$, con $u(\mathfrak{M}')\subset\mathfrak{N}'$,
$v(\mathfrak{N}')\subset\mathfrak{P}'$, de modo que se tienen homomorfismos
$u':A'\to B'$, $v':B'\to C'$, se obtienen homomorfismos canónicos
$u'_{C'/B'/A'}$, $v'_{C'/B'/A'}$ que, en virtud de
\sref{0.20.5.4-es}, forman sistemas proyectivos y dan por consiguiente,
pasando al límite, homomorfismos canónicos, prolongaciones a los completados
separados de los homomorfismos de la sucesión exacta
\sref{0.20.5.7.1-es}
\[
  \Omega^1_{B/A}\otimes_B C\xrightarrow{v_{C/B/A}}\Omega^1_{C/A}
  \xrightarrow{u_{C/B/A}}\Omega^1_{C/B}\to0:
\]
\begin{align}
  \widehat v_{C/B/A}:\widehat{\Omega}^1_{B/A}
  \widehat\otimes_{\widehat B}\widehat C&\to\widehat{\Omega}^1_{C/A},
  \tag{20.7.17.1}\label{0.20.7.17.1-es}\\
  \widehat u_{C/B/A}:\widehat{\Omega}^1_{C/A}&\to
  \widehat{\Omega}^1_{C/B}. \tag{20.7.17.2}\label{0.20.7.17.2-es}
\end{align}
y en la sucesión
\begin{equation}
\tag{20.7.17.3}
\label{0.20.7.17.3-es}
  \widehat{\Omega}^1_{B/A}\widehat\otimes_{\widehat B}\widehat C
  \xrightarrow{\widehat v_{C/B/A}}\widehat{\Omega}^1_{C/A}
  \xrightarrow{\widehat u_{C/B/A}}\widehat{\Omega}^1_{C/B}\to0
\end{equation}
la composición de dos homomorfismos consecutivos es $0$, pero la sucesión
\emph{no es necesariamente exacta}. Sin embargo, si $B$ y $C$ son
\emph{metrizable}, el homomorfismo $\widehat u_{C/B/A}$ es \emph{sobreyectivo},
y $\operatorname{Im}(\widehat v_{C/B/A})$ es \emph{densa} en
$\operatorname{Ker}(\widehat u_{C/B/A})$: esto se deduce inmediatamente
(cf.\ ($\mathbf{0}_{\mathrm I}$, 7.3.1)) del hecho de que, si
$(\mathfrak{M}_k)$ (resp.\ $(\mathfrak{N}_k)$) es una sucesión decreciente
de ideales de $B$ (resp.\ $C$), que forma un sistema fundamental de entornos
de $0$, y si se ponen $B_k=B/\mathfrak{M}_k$, $C_k=C/\mathfrak{N}_k$, los
homomorfismos de transición
$\Omega^1_{C_{k+1}/B_{k+1}}\to\Omega^1_{C_k/B_k}$,
$\Omega^1_{C_{k+1}/A}\to\Omega^1_{C_k/A}$ y
$\Omega^1_{B_{k+1}/A}\to\Omega^1_{B_k/A}$ son sobreyectivos
\sref{0.20.7.14-es}.
\end{env}

\begin{proposition}[20.7.18]
\label{0.20.7.18-es}
Sean $A$, $B$, $C$ tres anillos topológicos, y sean $u:A\to B$, $v:B\to C$
dos homomorfismos continuos. Se supone que $B$ y $C$ son admisibles
($\mathbf{0}_{\mathrm I}$, 7.1.2) y metrizables. Para que el homomorfismo
canónico
\[
  \widehat v_{C/B/A}:\widehat{\Omega}^1_{B/A}
  \widehat\otimes_{\widehat B}\widehat C\to\widehat{\Omega}^1_{C/A}
\]
admita una inversa por la izquierda que sea un $C$-homomorfismo continuo, es
necesario y suficiente que $C$ sea una $B$-álgebra formalmente lisa respecto
de $A$.
\end{proposition}

\begin{proof}
La condición es necesaria en virtud de \sref{0.20.7.2-es} y
\sref{0.19.1.6-es}. Para ver que es suficiente, obsérvese que el $C$-módulo
topológico
$L=\widehat{\Omega}^1_{B/A}\widehat\otimes_{\widehat B}\widehat C$ es
metrizable y
\oldpage[0\textsubscript{IV-1}]{153}
completo en virtud de la hipótesis, por lo que
$E=D_C(L)$, provisto de la topología producto, es metrizable y completo; es
además admisible, pues si $\mathfrak{R}$ es un ideal de definición en $C$, la
sucesión $(\mathfrak{R}\oplus L)^n=\mathfrak{R}^n\oplus\mathfrak{R}^{n-1}L$
tiende a $0$. Como la aplicación compuesta
\[
  D:B\xrightarrow{d_{B/A}}\Omega^1_{B/A}\to
  \widehat{\Omega}^1_{B/A}\widehat\otimes_{\widehat B}\widehat C=L
\]
es una $A$-derivación continua de $B$ en $L$, el $A$-homomorfismo continuo
\[
  f:x\longmapsto(v(x),D(x))
\]
de $B$ en $E$ define en $E$ una estructura de extensión de $B$. Como $L$ es
un ideal cerrado en $E$, se deduce de \sref{0.19.9.5-es} y de la hipótesis que
la aplicación identidad $C\to E/L$ (que es un $B$-homomorfismo) se factoriza
como $C\xrightarrow{g}E\to E/L$, donde $g$ es un homomorfismo continuo tal
que $g\circ v=f$; por consiguiente, según \sref{0.20.1.3-es}, $g$ es de la
forma $y\mapsto(y,D'(y))$, donde $D'$ es una $B$-derivación continua de $C$ en
$L$, es decir, $D'\circ v=D$. Teniendo en cuenta \sref{0.20.7.14.4-es}, se
tiene $D'=h\circ\widehat d_{C/A}$, donde
$h:\widehat{\Omega}^1_{C/A}\to L$ es un $C$-homomorfismo continuo; pero por
definición, $\widehat d_{C/A}\circ v=\widehat v_{C/B/A}\circ D$, y como la
imagen de $B$ por $D$ genera (topológicamente) el $C$-módulo $L$
\sref{0.20.4.7-es}, la relación $D'\circ v=D$ da precisamente
$h\circ\widehat v_{C/B/A}=1_L$.
\end{proof}

\begin{corollary}[20.7.19]
\label{0.20.7.19-es}
Bajo las hipótesis de \sref{0.20.7.18-es}, supongamos además que en $B$ y
$C$ el cuadrado de todo ideal abierto sea abierto; entonces, para que $C$ sea
una $B$-álgebra formalmente lisa respecto de $A$, es necesario y suficiente
que $\widehat v_{C/B/A}$ sea invertible por la izquierda.
\end{corollary}

\begin{proof}
En efecto, las topologías de
$\widehat{\Omega}^1_{B/A}\widehat\otimes_{\widehat B}\widehat C$ y de
$\widehat{\Omega}^1_{C/A}$ se deducen entonces de la de $C$
\sref{0.20.4.5-es}, y todo $C$-homomorfismo de uno en el otro es
necesariamente continuo.
\end{proof}

\begin{env}[20.7.20]
\label{0.20.7.20-es}
Sean $A$ un anillo topológico, $B$ una $A$-álgebra topológica
\emph{metrizable y completa}, $\mathfrak{R}$ un ideal \emph{cerrado} de $B$,
$C=B/\mathfrak{R}$ el anillo topológico cociente, que es metrizable y
completo. Sea $(\mathfrak{M}_k)$ un sistema fundamental decreciente de
entornos de $0$ en $B$ formado por ideales, y pongamos
$B_k=B/\mathfrak{M}_k$,
$\mathfrak{R}_k=(\mathfrak{R}+\mathfrak{M}_k)/\mathfrak{M}_k$,
$C_k=B_k/\mathfrak{R}_k$. Se tiene un sistema proyectivo de homomorfismos
\[
  \delta_{C_k/B_k/A}:\mathfrak{R}_k/\mathfrak{R}_k^2\to
  \Omega^1_{B_k/A}\otimes_{B_k}C_k
\]
\sref{0.20.5.11.3-es}, de donde se deduce, pasando al límite, un
homomorfismo canónico
\[
  \widehat\delta_{C/B/A}:\mathfrak{R}/\mathfrak{R}^2\to
  \widehat{\Omega}^1_{B/A}\widehat\otimes_{\widehat B}\widehat C
\]
y, razonando como en \sref{0.20.7.17-es}, se ve que el homomorfismo canónico
$\widehat u_{C/B/A}:\widehat{\Omega}^1_{B/A}
\widehat\otimes_{\widehat B}\widehat C\to\widehat{\Omega}^1_{C/A}$ es
\emph{sobreyectivo} y que $\operatorname{Im}(\widehat\delta_{C/B/A})$ es
\emph{densa} en $\operatorname{Ker}(\widehat u_{C/B/A})$.
\end{env}
